%% ============================================================
%%  MASTER TEX FILE
%%  Applications of Generative Orthogonal Matrix Compression Science
%%  The Complete Principia Orthogona Series, Volumes I-III with Applications
%%  Author: Pablo Nogueira Grossi
%%  Publisher: G6 LLC, Newark, NJ, USA
%%  ISBN (Print):  979-8-9954416-0-1
%%  ISBN (eBook):  979-8-9954416-1-8
%%  Year: 2026
%% ============================================================

\documentclass[12pt, oneside]{book}

%% ── PACKAGES ────────────────────────────────────────────────
\usepackage[T1]{fontenc}
\usepackage[utf8]{inputenc}
\usepackage{amsmath,amssymb,amsthm,amsfonts}
\usepackage{mathrsfs}
\usepackage{mathtools}
\usepackage{geometry}
\usepackage{enumitem}
\usepackage{microtype}
\usepackage{booktabs}
\usepackage{array}
\usepackage{multirow}
\usepackage{emptypage}
\usepackage{fancyhdr}
\usepackage[hidelinks]{hyperref}

\geometry{
  paperwidth=5.5in,
  paperheight=8.5in,
  top=0.875in,
  bottom=0.875in,
  inner=0.875in,
  outer=0.625in,
  headheight=14pt
}

%% ── PAGE STYLE ───────────────────────────────────────────────
\pagestyle{fancy}
\fancyhf{}
\fancyhead[LE]{\small\itshape Principia Orthogona}
\fancyhead[RO]{\small\itshape Pablo Nogueira Grossi}
\fancyfoot[C]{\small\thepage}
\renewcommand{\headrulewidth}{0.4pt}
\fancypagestyle{plain}{
  \fancyhf{}
  \fancyfoot[C]{\small\thepage}
  \renewcommand{\headrulewidth}{0pt}
}

%% ── THEOREM ENVIRONMENTS ─────────────────────────────────────
\theoremstyle{plain}
\newtheorem{theorem}{Theorem}[section]
\newtheorem{lemma}[theorem]{Lemma}
\newtheorem{corollary}[theorem]{Corollary}
\newtheorem{proposition}[theorem]{Proposition}
\newtheorem{conjecture}[theorem]{Conjecture}
\theoremstyle{definition}
\newtheorem{definition}{Definition}[section]
\newtheorem{assumption}{Assumption}[section]
\newtheorem{axiom}{Axiom}
\newtheorem{invariant}{Invariant}[section]
\newtheorem{normalform}{Normal Form}[section]
\theoremstyle{remark}
\newtheorem{remark}{Remark}[section]
\newtheorem{example}{Example}[section]

%% ── MACROS (consolidated from all four papers) ───────────────
\newcommand{\R}{\mathbb{R}}
\newcommand{\T}{\mathbb{T}}
\newcommand{\dm}{\mathfrak{D}}
\newcommand{\dmcat}{\mathbf{dm3}}
\newcommand{\Orb}{\Gamma}
\newcommand{\cN}{\mathcal{N}}
\newcommand{\cB}{\mathcal{B}}
\newcommand{\cS}{\mathcal{S}}
\newcommand{\cL}{\mathcal{L}}
\newcommand{\cR}{\mathcal{R}}
\newcommand{\cV}{\mathcal{V}}
\newcommand{\cM}{\mathcal{M}}
\newcommand{\cU}{\mathcal{U}}
\newcommand{\cP}{\mathcal{P}}
\newcommand{\cD}{\mathcal{D}}
\newcommand{\eps}{\varepsilon}
\newcommand{\dg}{d_g}
\newcommand{\norm}[1]{\left\|#1\right\|}
\newcommand{\abs}[1]{\left|#1\right|}
\DeclareMathOperator{\lcm}{lcm}
\DeclareMathOperator{\Fix}{Fix}
\DeclareMathOperator{\Hess}{Hess}
\DeclareMathOperator{\supp}{supp}
\DeclareMathOperator{\rank}{rank}

%% ── HYPERREF SETUP ───────────────────────────────────────────
\hypersetup{
  pdftitle={Applications of Generative Orthogonal Matrix Compression Science},
  pdfauthor={Pablo Nogueira Grossi},
  pdfsubject={Mathematics, Contact Geometry, Biological Systems},
  pdfkeywords={contact geometry, generative transitions, dm3, Principia Orthogona, G6 LLC},
  pdfcreator={G6 LLC},
  colorlinks=false
}

%% ════════════════════════════════════════════════════════════
\begin{document}
%% ════════════════════════════════════════════════════════════

%% ── FRONT MATTER (Roman numerals) ────────────────────────────
\frontmatter

%% Half title
\thispagestyle{empty}
\vspace*{4cm}
\begin{center}
  {\large Applications of}\\[0.5em]
  {\LARGE\bfseries Generative Orthogonal\\[0.3em]
  Matrix Compression Science}
\end{center}
\cleardoublepage

%% Title page
\thispagestyle{empty}
\begin{center}
  \vspace*{1.5cm}
  {\LARGE\bfseries Applications of}\\[0.4em]
  {\Large\bfseries Generative Orthogonal\\[0.3em]
  Matrix Compression Science}\\[1.5em]
  \rule{0.6\textwidth}{0.8pt}\\[1em]
  {\large\itshape The Complete Principia Orthogona Series\\
  Volumes I--III with Applications}\\[2em]
  \rule{0.6\textwidth}{0.4pt}\\[2em]
  {\large\bfseries Pablo Nogueira Grossi}\\[0.5em]
  {\normalsize G6 LLC}\\[4em]
  \vfill
  \rule{0.6\textwidth}{0.8pt}\\[0.5em]
  {\small G6 LLC \textperiodcentered\ Newark, New Jersey \textperiodcentered\ 2026}
\end{center}
\cleardoublepage

%% Copyright page
\thispagestyle{empty}
\vspace*{2cm}
{\small
\noindent\textbf{Applications of Generative Orthogonal Matrix Compression Science}\\
\textit{The Complete Principia Orthogona Series, Volumes I--III with Applications}

\bigskip
\noindent Copyright \copyright\ 2025--2026 Pablo Nogueira Grossi / G6 LLC.\\
All rights reserved.

\bigskip
\noindent No part of this publication may be reproduced, distributed, or transmitted
in any form or by any means without the prior written permission of the publisher,
except brief quotations embodied in scholarly reviews.

\bigskip
\noindent Author: Pablo Nogueira Grossi\\
Publisher: G6 LLC, Newark, NJ, USA\\
Imprint: G6 LLC\\
Contact: pablogrossi@hotmail.com\\
ORCID: 0009-0000-6496-2186

\bigskip
\noindent Preprint: Zenodo DOI 10.5281/zenodo.19117400\\
HAL open science: hal-05555216, hal-05559997

\bigskip
\noindent\textbf{Print ISBN: 979-8-9954416-0-1}\\
\textbf{eBook ISBN: 979-8-9954416-1-8}

\bigskip
\noindent First edition, 2026\\
Printed in the United States of America

\bigskip
{\footnotesize All claims in the biological applications sections are subject to
further empirical research and peer review.}
}
\cleardoublepage

%% Dedication
\thispagestyle{empty}
\vspace*{6cm}
\begin{center}
  \itshape
  Dedicated to my children\\[0.5em]
  Vic (R.I.P.), Giulia, Alice (R.I.P.), Sarah (R.I.P.), and David.\\[1em]
  \normalfont\textit{Once tiny, always strong.}
\end{center}
\cleardoublepage

%% Series Foreword
\chapter*{Series Foreword}
\addcontentsline{toc}{chapter}{Series Foreword}

This volume collects the complete Principia Orthogona series, a unified mathematical
framework for generative transitions developed independently by Pablo Nogueira Grossi
over a period of more than twenty-five years, beginning in the year 2000 and completed
in draft form by 2024.

The series develops a single governing idea across three layers:

\textbf{Volume I --- The Mathematics of Generative Transitions} introduces the operator
sequence $C \longrightarrow K \longrightarrow F \longrightarrow U$ on a Riemannian
manifold, governing compression, curvature intensification, fold (loss of injectivity),
and stabilization. This volume establishes the geometric precursor of all subsequent
structure: the intrinsic curvature threshold $\kappa^*$, the Whitney $A_1$--$A_3$
singularity classification, and the free-discontinuity variational principle.

\textbf{Volume II --- Generative Contact Mechanics (dm3)} constructs the
contact-geometric realization of Volume I. The central object is the dm3
system---a smooth Riemannian manifold with a hyperbolic limit cycle, Lyapunov function,
and stochastic extension, formalized through eight axioms. Four main theorems establish
existence and stability of limit cycles, closure under a unification operator, a
universal contact normal form, and $C^1$-structural stability with an explicit radius
$\varepsilon_0 = 1/3$.

\textbf{Volume III --- Biological Instantiations} shows that the dm3 framework is
instantiated directly in four classes of oscillatory biological systems: the HPA
allostatic stress network, neural oscillations, circadian rhythms, and immune
adaptation cycles. Three falsifiable quantitative predictions follow from the
abstract theorems.

\textbf{Applications of GOMC Science} presents the cross-volume synthesis and
the broader applications of Generative Orthogonal Matrix Compression Science.

The mathematical portions are typeset with full theorem-proof structure, MSC 2020
subject classifications, and bibliographies of peer-reviewed references. Preprint
versions are deposited on Zenodo and HAL under open-access licenses.

\bigskip
\begin{flushright}
Pablo Nogueira Grossi\\
Newark, New Jersey\\
March 2026
\end{flushright}

\cleardoublepage

%% Table of Contents
\tableofcontents
\cleardoublepage

%% ── MAIN MATTER (Arabic numerals) ────────────────────────────
\mainmatter


%% ════════════════════════════════════════════════════════════
\part{The Mathematics of Generative Transitions}
%% ════════════════════════════════════════════════════════════

\begin{center}
\emph{Dedicated to my children Vic~(R.I.P.), Giulia,
  Alice~(R.I.P.), Sarah~(R.I.P.), and David.}\\[4pt]
\emph{Once tiny, always strong.}
\end{center}

\vspace{1.5em}

\newpage

%-----------------------------------------------------------------------
\section*{Preface}
%-----------------------------------------------------------------------

This volume develops a unified mathematical framework for generative
transitions: localized geometric events in which a trajectory undergoes
compression, curvature intensification, loss of injectivity, and
stabilization.

The central object is the operator sequence
\[
C \;\rightarrow\; K \;\rightarrow\; F \;\rightarrow\; U,
\]
acting on trajectories in a Riemannian manifold.

Several results are presented with proof sketches rather than complete
proofs. These include the symplectic preservation across the fold, the
classification theorem for admissible singularities, the existence and
uniqueness of the unfolding map, and the stability of $\kappa^*$ under
perturbations. Each can be completed with standard tools from
differential geometry, singularity theory, and variational analysis;
full proofs lie beyond the scope of this foundational volume.

This volume does not claim that all physical, biological, cognitive,
or economic systems instantiate this structure. It provides the
mathematical substrate on which such questions can be posed rigorously.
Cross-domain instantiation is the subject of Volume Three.

This volume is the first in a three-part series. The subsequent
papers---\emph{Generative Contact Mechanics: A Geometric Framework for
Dissipative Systems with Structured Limit Cycles} and \emph{The dm3
Operator: Explicit Toy Model and Global Dynamical Analysis}---develop
the full contact-geometric realization of the generative transition
theory introduced here.

\medskip

\noindent\textbf{Note to the reader.}
Readers who prefer to begin with a concrete realization before
engaging with the abstract framework are invited to consult
Volume~Two first, which constructs an explicit two-dimensional
contact-Hamiltonian system instantiating every definition, operator,
and theorem of this volume. In the words of that paper:
\emph{the dm3 Toy Model is a complete verification witness for
Generative Contact Mechanics, demonstrating that the entire abstract
framework is jointly realizable, computable, and dynamically consistent.}
The present volume is logically prior; Volume~Two is pedagogically
illuminating as a first reading.

%-----------------------------------------------------------------------
\section{Overview}
%-----------------------------------------------------------------------

Let $(X, g)$ be a smooth Riemannian manifold. A \emph{generative
transition} is a localized event along a trajectory
$\gamma : [0,T] \to X$ consisting of four sequential phases:
compression, curvature induction, folding, and stabilization.

The theory is organized around the composite operator
\[
\mathcal{G} = U \circ F \circ K \circ C : X \to X.
\]

The manifold $N = (M, \Phi, F, K, g)$ serves as the canonical
instance, where $M$ is the geometric substrate, $\Phi$ the stability
functional, $F$ the directional field, $K$ the constraint operator,
and $g$ the generative operator.

The generative flow on $M$ is governed by
\[
\frac{dx}{dt} = -\nabla \Phi(x) + F(x) + \eta(t),
\]
where $\eta(t)$ is a stochastic perturbation.

%-----------------------------------------------------------------------
\section{Minimal Mathematical Assumptions}
%-----------------------------------------------------------------------

The operator sequence is defined under the weakest assumptions
necessary for the constructions of Section~\ref{sec:operators} to be
well-posed.

\begin{assumption}[Manifold Structure]\label{ass:manifold}
The state space $X$ is a smooth, finite-dimensional Riemannian
manifold, locally compact and second-countable.
\end{assumption}

\begin{assumption}[Trajectory Regularity]\label{ass:traj}
A system trajectory $\gamma : [0,T] \to X$ is piecewise $C^2$,
locally non-degenerate ($\|\dot\gamma(t)\| \neq 0$ a.e.), and has
bounded curvature on compact intervals prior to folding events.
\end{assumption}

\begin{assumption}[Compression Feasibility]\label{ass:compress}
There exists a lower-dimensional submanifold $X_C \subset X$ and a
Lipschitz continuous projection $C : X \to X_C$ satisfying the
non-collapse condition
\[
d\bigl(C(x_1), C(x_2)\bigr) \;\ge\; \delta\, d(x_1, x_2)
\]
for some $\delta > 0$ and all $x_1, x_2$ in a compact neighborhood.
\end{assumption}

\begin{remark}
Assumption~\ref{ass:compress} is a bi-Lipschitz embedding condition
on the compression map, ensuring distinct trajectories remain
distinguishable after compression. This is stronger than mere Lipschitz
continuity and should be understood as a lower bound on distance
preservation.
\end{remark}

\begin{assumption}[Curvature Threshold]\label{ass:kappa}
The critical curvature is defined intrinsically by the focal radius:
\[
\kappa^*(x) \;=\; \frac{1}{\operatorname{foc}(x)},
\]
where $\operatorname{foc}(x)$ is the distance from $x$ to the first
focal point along the normal direction. In the presence of positive
sectional curvature $K_{\mathrm{sec}}(x) > 0$, this is modified by
the Rauch comparison theorem to
\[
\kappa^*(x) \;=\; \min\!\left(\|II_x\|,\;
  \sqrt{K_{\mathrm{sec}}(x)}\right),
\]
where $\|II_x\|$ is the norm of the second fundamental form.
For $K_{\mathrm{sec}}(x) \le 0$, we have $\kappa^*(x) = \|II_x\|$.
\end{assumption}

\begin{assumption}[Folding Well-Posedness]\label{ass:fold}
The folding operator $F : X_C \to X_F$ satisfies: the Jacobian $dF$
loses rank by exactly $1$ at fold points; the fold is local; and the
fold produces a finite number of branches.
\end{assumption}

\begin{assumption}[Morse Stability Functional]\label{ass:morse}
The stability functional $\Phi : X \to \mathbb{R}$ is $C^2$, bounded
below on compact subsets, and \emph{Morse}: all critical points are
non-degenerate, i.e., $\nabla^2\Phi(x^*) \succ 0$ at every local
minimum~$x^*$.
\end{assumption}

%-----------------------------------------------------------------------
\section{Operator Definitions}\label{sec:operators}
%-----------------------------------------------------------------------

\subsection{State Space}

Let $X$ be as in Assumption~\ref{ass:manifold}. A system trajectory
is a piecewise $C^2$ map $\gamma : [0,T] \to X$. The goal is to
describe local transitions in $\gamma$ via the sequence
$C \to K \to F \to U$.

\subsection{Compression Operator \texorpdfstring{$C$}{C}}

\begin{definition}[Compression]
A \emph{compression operator} is a map $C : X \to X_C$ with
$\dim(X_C) < \dim(X)$ satisfying Assumption~\ref{ass:compress}.
\end{definition}

$C$ reduces degrees of freedom while retaining essential local
structure.

\subsection{Curvature Operator \texorpdfstring{$K$}{K}}

\begin{definition}[Curvature Operator]
Given a compressed trajectory $\gamma_C : [0,T] \to X_C$, the
curvature operator $K : X_C \to X_C$ modifies the tangent field by
\[
\frac{d}{ds}\bigl(K(\gamma_C)(s)\bigr)
\;=\; \dot\gamma_C(s) + \alpha(s)\,n(s),
\]
where $n(s)$ is the unit normal and
\[
\alpha(s) \;=\; \lambda\,\bigl(\kappa^*(\gamma_C(s))
  - \kappa(s)\bigr)_+, \quad \lambda > 0.
\]
\end{definition}

$K$ is a curvature-inducing flow that drives curvature monotonically
toward $\kappa^*$ but never beyond it. Consequently, $K$ alone does
not trigger folding; the sequence is a \emph{hybrid dynamical system}
with a switching condition at $\kappa^*$.

\subsection{Folding Operator \texorpdfstring{$F$}{F}}

\begin{definition}[Folding Map]
Let $\gamma_K = K(\gamma_C)$. A \emph{fold} occurs at $s_0$ when
$|\kappa_K(s_0)| = \kappa^*(\gamma_K(s_0))$. The folding operator
$F : X_C \to X_F$ acts by
\[
F(\gamma_K(s)) \;=\; \gamma_K(s) - \beta(s)\,n(s),
\]
where
\[
\beta(s) \;=\;
\begin{cases}
0, & |\kappa_K(s)| < \kappa^*, \\
\mu\bigl(|\kappa_K(s)| - \kappa^*\bigr), & |\kappa_K(s)| \ge \kappa^*.
\end{cases}
\]
Here $\mu > 0$ controls fold depth.
\end{definition}

At a fold point $s_0$,
$\operatorname{rank}(dF_{\gamma_K(s_0)}) = \dim(X_C) - 1$.

\subsection{Unfolding Operator \texorpdfstring{$U$}{U}}

\begin{definition}[Unfolding Map]
The unfolding operator $U : X_F \to X$ is
\[
U(x_F) \;=\; \operatorname*{arg\,min}_{y \in \mathcal{N}(x_F)}
  \Phi(y),
\]
where $\mathcal{N}(x_F)$ is a sufficiently small neighborhood of
$x_F$ and $\Phi$ satisfies Assumption~\ref{ass:morse}.
\end{definition}

The unfolding process is realised by the gradient flow
$\dot y = -\nabla\Phi(y)$, $y(0) = x_F$, which converges to a
non-degenerate local minimum~$x^*$.

\subsection{Compatibility of \texorpdfstring{$K$}{K} and
\texorpdfstring{$F$}{F}}

\begin{theorem}[Sequential Consistency]
If $K$ is applied until curvature reaches $\kappa^*$, then $F$ is
well-defined, produces a finite branch set, and induces a
rank-deficient Jacobian at fold points.
\end{theorem}

\begin{proof}[Proof sketch]
$K$ drives curvature monotonically to $\kappa^*$ via the gain field
$\alpha(s)$. At $\kappa^*$, the term $\beta(s)$ becomes nonzero,
introducing local non-injectivity. The normal direction collapses,
reducing Jacobian rank by~$1$. Locality and the finite critical-point
structure of $\kappa^*$ (via the Morse assumption on $\Phi$) imply
finitely many branches.
\end{proof}

%-----------------------------------------------------------------------
\section{Falsifiability Conditions}
%-----------------------------------------------------------------------

The framework is falsifiable at four levels.

\paragraph{Compression.}
If empirical data show expansion under $C$, or collapse of distinct
trajectories, the model fails.

\paragraph{Curvature.}
If a fold occurs at curvature strictly below $\kappa^*$, or curvature
exceeds $\kappa^*$ without folding, the model is falsified. ($\kappa^*$
is computed independently via the focal radius, not post-hoc from fold
observations.)

\paragraph{Folding.}
If empirical fold events do not correspond to Jacobian rank loss, or
produce infinitely many branches, the model fails.

\paragraph{Sequence order.}
If transitions occur in a different order, or stabilization occurs
without folding, the model is falsified.

%-----------------------------------------------------------------------
\section{Structural Theorems}
%-----------------------------------------------------------------------

\begin{theorem}[Existence and Well-Posedness]
Under Assumptions~\ref{ass:manifold}--\ref{ass:morse}, the composite
operator $\mathcal{G} = U \circ F \circ K \circ C$ is well-defined on
any piecewise $C^2$ trajectory.
\end{theorem}

\begin{theorem}[Local Determination]
The action of $\mathcal{G}$ on $\gamma$ is determined entirely by the
local geometry of $\gamma$ in a neighborhood of the fold point.
\end{theorem}

\begin{theorem}[Non-Commutativity]
The operators $C$, $K$, $F$, $U$ do not commute; the sequence is
order-dependent.
\end{theorem}

\begin{theorem}[Irreducibility]
No operator in the sequence $C \to K \to F \to U$ can be removed
without altering the qualitative structure of the transition.
\end{theorem}

\begin{theorem}[Finite Branching]
The branch set $B = \{F(\gamma_K(s_i)) :
|\kappa_K(s_i)| = \kappa^*(\gamma_K(s_i))\}$ is finite.
\end{theorem}

\begin{remark}
Finiteness follows from the Morse condition on $\Phi$
(Assumption~\ref{ass:morse}), which prevents fourth-order degeneracies
and ensures isolated critical points, together with the transversality
of $\gamma$ to the fold locus (a generic condition to be assumed
explicitly in applications).
\end{remark}

%-----------------------------------------------------------------------
\section{Canonical Examples}
%-----------------------------------------------------------------------

\paragraph{Planar curve with curvature-driven flow.}
The simplest nontrivial system: $\gamma : [0,T] \to \mathbb{R}^2$,
with $C$ an orthogonal projection, $K$ the curvature-inducing flow,
$F$ activated at $\kappa^*$, and $U$ gradient descent on a local
$\Phi$.

\paragraph{Elastic rod under compression.}
An elastic rod minimising Euler--Bernoulli energy
$E[\gamma] = \int \kappa(s)^2\,ds$. Axial loading provides $C$;
buckling provides $K$; the critical load is $\kappa^*$; post-buckling
configuration provides $U$. This is the cleanest physical realisation
of the curvature threshold.

\paragraph{Gradient flow on a double-well potential.}
$\Phi(x) = (x^2 - 1)^2$. The unstable equilibrium at $x = 0$ is the
fold point; gradient flow selects one of $x = \pm 1$. Illustrates
finite branching and stability selection.

\paragraph{Saddle-node bifurcation.}
$\dot x = \mu - x^2$, $\dot y = -y$. Projection onto the slow
manifold provides $C$; approach to the fold point provides $K$; loss
of the slow manifold at $\mu = 0$ provides $F$; flow to the stable
branch provides $U$.

\paragraph{Plasma sheet and market manifolds (Volume Three).}
Magnetic field-line reconnection and price-trajectory bifurcations are
conjectured instantiations. Each requires domain-specific construction
of $(X, g, \Phi)$ and empirical computation of $\kappa^*$. These are
deferred to Volume Three.

%-----------------------------------------------------------------------
\section{Analytical Invariants}
%-----------------------------------------------------------------------

\begin{invariant}[Ambient Dimension]
$\dim(X)$ is preserved by $\mathcal{G}$. No claim is made about the
dimension of the image of a specific trajectory.
\end{invariant}

\begin{invariant}[Codimension of the Fold]
$\operatorname{codim}(F(\gamma)) = 1$, following from the rank-loss
condition.
\end{invariant}

\begin{invariant}[Critical Curvature Threshold]
$\kappa^*(x)$ is a geometric property of the manifold, invariant under
the operator sequence.
\end{invariant}

\begin{invariant}[Curvature Sign]
$\operatorname{sgn}(\kappa(s))$ is preserved under $K$ and $F$; only
magnitude changes.
\end{invariant}

\begin{invariant}[Injectivity Before Threshold]
For all $s$ with $|\kappa(s)| < \kappa^*(s)$, the trajectory remains
injective under $C$, $K$, $F$.
\end{invariant}

\begin{invariant}[Rank Deficiency at Fold]
$\operatorname{rank}(dF) = \dim(X_C) - 1$ at every fold point.
\end{invariant}

\begin{invariant}[Energy Monotonicity under $U$]
$\Phi(U(x)) \le \Phi(x)$, with strict inequality unless $x$ is
already a local minimum.
\end{invariant}

\begin{remark}
An energy-gap invariant $\Delta\Phi = \Phi(x_2) - \Phi(x_1)$ is
preserved under $C$, $K$, $F$ only if $\Phi$ is constant along the
fibers of each operator. This is not guaranteed in general and is not
claimed as a universal invariant.
\end{remark}

%-----------------------------------------------------------------------
\section{Normal Forms}
%-----------------------------------------------------------------------

Two transitions $\mathcal{G}_1$, $\mathcal{G}_2$ are equivalent if
there exists a diffeomorphism $\psi : X \to X$ with
$\mathcal{G}_2 = \psi \circ \mathcal{G}_1 \circ \psi^{-1}$.

\begin{normalform}[Compression]
$C_{\mathrm{NF}} = \pi_k : \mathbb{R}^n \to \mathbb{R}^k$, the
standard coordinate projection.
\end{normalform}

\begin{normalform}[Curvature]
$K_{\mathrm{NF}}(s) = (s,\, \alpha s^2)$, $\alpha > 0$. This
produces curvature
$\kappa(s) = 2\alpha(1 + 4\alpha^2 s^2)^{-3/2}$, which approximates
$2\alpha$ near $s = 0$.
\end{normalform}

\begin{normalform}[Folding --- Whitney Fold]
$F_{\mathrm{NF}}(u,v) = (u, v^2)$. This is the unique (up to
diffeomorphism) normal form for a codimension-1 fold singularity.
\end{normalform}

\begin{normalform}[Unfolding]
$U_{\mathrm{NF}}(x) = 0$, the canonical stabilization map arising
from $\Phi_{\mathrm{NF}}(x) = x^2$.
\end{normalform}

%-----------------------------------------------------------------------
\section{Singularity Classification}
%-----------------------------------------------------------------------

\subsection{Admissible Singularities}

Given the constraints of rank-1 loss, finite branching, and the Morse
condition on $\Phi$, the admissible singularities are precisely:

\begin{enumerate}
\item Fold $A_1$: $(u,v) \mapsto (u, v^2)$
\item Cusp $A_2$: $(u,v) \mapsto (u, v^3 + uv)$
\item Swallowtail $A_3$: $(u,v) \mapsto (u, v^4 + uv^2 + \beta v)$
\end{enumerate}

Higher $A_k$ ($k \ge 4$) require a $k$-parameter unfolding. The Morse
condition on $\Phi$ restricts the unfolding to at most three
parameters, excluding $A_4$ and above.

\subsection{Classification Conditions}

Let $\Delta(s) = \kappa(s) - \kappa^*(s)$.

\begin{center}
\begin{tabular}{lll}
\toprule
Type & Conditions & Normal form \\
\midrule
$A_1$ (fold) & $\Delta(s_0)=0$, $\Delta'(s_0)\ne 0$ & $(u,v^2)$ \\
$A_2$ (cusp) & $\Delta=\Delta'=0$, $\Delta''\ne 0$ & $(u,v^3+uv)$ \\
$A_3$ (swallowtail) & $\Delta=\Delta'=\Delta''=0$, $\Delta'''\ne 0$
  & $(u,v^4+uv^2+\beta v)$ \\
\bottomrule
\end{tabular}
\end{center}

\begin{theorem}[Classification of Generative Transitions]
Every admissible generative transition $\mathcal{G} = U \circ F \circ
K \circ C$ is $\mathcal{A}$-equivalent to exactly one of $A_1$,
$A_2$, $A_3$.
\end{theorem}

\begin{proof}[Proof sketch]
Rank loss is exactly $1$, restricting to the $A_k$ series. The Morse
condition on $\Phi$ limits the unfolding to at most three parameters.
Transversality of $\gamma$ to the fold locus (generic assumption)
ensures isolated fold points. Thus $k \le 3$.
\end{proof}

\subsection{Bifurcation Stratification}

Parameter space $\Theta \subset \mathbb{R}^p$, $p \le 3$, decomposes
into:
\begin{itemize}
\item $\Theta_1$: generic fold region (codimension 0 in $\Theta$),
\item $\Theta_2$: cusp stratum (codimension 1),
\item $\Theta_3$: swallowtail stratum
  (codimension 2; a single point when $p=3$).
\end{itemize}

%-----------------------------------------------------------------------
\section{Metric Geometry of \texorpdfstring{$\kappa^*$}{kappa*}}
\label{sec:metric}
%-----------------------------------------------------------------------

\subsection{Intrinsic Definition}

\[
\kappa^*(x) \;=\; \frac{1}{\operatorname{foc}(x)},
\]
where $\operatorname{foc}(x)$ is the focal radius. For a curve
embedded in $(X,g)$, this equals $\|II_x\|$ (the norm of the second
fundamental form) when $K_{\mathrm{sec}} \le 0$.

\subsection{Sectional Curvature Correction}

By the Rauch comparison theorem, for variable positive sectional
curvature:
\[
\kappa^*(x) \;=\; \min\!\left(\|II_x\|,\;
  \sqrt{K_{\mathrm{sec}}(x)}\right).
\]
Positive ambient curvature lowers the threshold for folding.

\subsection{Computability}

Given $\gamma$ in $(X,g)$:
\begin{enumerate}
\item Compute $\|II_x\|$ from the second fundamental form.
\item Compute $K_{\mathrm{sec}}(x)$ from the curvature tensor.
\item Apply the formula above.
\end{enumerate}

$\kappa^*$ is stable:
$\delta\kappa^* = O(\delta g) + O(\delta II)
  + O(\delta K_{\mathrm{sec}})$.

%-----------------------------------------------------------------------
\section{Variational Principles}
%-----------------------------------------------------------------------

Define the action functional
\[
\mathcal{S}[\gamma] \;=\; \int_0^T
\Bigl(
\tfrac{1}{2}\|P_\perp \dot\gamma\|^2
+ \tfrac{\lambda}{2}(\kappa^* - \kappa)_+^2
+ \mu\,\delta(|\kappa| - \kappa^*)
+ \Phi(\gamma)
\Bigr)\,ds.
\]

Each term corresponds to one operator:
\begin{itemize}
\item $L_C = \tfrac{1}{2}\|P_\perp\dot\gamma\|^2$: minimise
  orthogonal kinetic energy (compression).
\item $L_K = \tfrac{\lambda}{2}(\kappa^*-\kappa)_+^2$: minimise
  curvature deficit.
\item $L_F = \mu\,\delta(|\kappa|-\kappa^*)$: singular activation
  at threshold.
\item $L_U = \Phi(\gamma)$: minimise stability potential.
\end{itemize}

The delta term places the framework in the class of
\emph{free-discontinuity variational problems}
(cf.\ Ambrosio--Tortorelli, Mumford--Shah, Francfort--Marigo).
It produces the jump condition
\[
\left[\frac{\partial L}{\partial \dot\gamma}\right]_{s_0^-}^{s_0^+}
= \mu\,n(s_0),
\]
which is the variational counterpart of the fold map.

The full generative transition is therefore a
\emph{piecewise-smooth extremal} of $\mathcal{S}$.

%-----------------------------------------------------------------------
\section{Hamiltonian Discontinuities and Symplectic Geometry}
%-----------------------------------------------------------------------

\subsection{Phase Space}

The phase space is $T^*X$ with canonical coordinates $(\gamma, p)$
and symplectic form $\omega = d\gamma \wedge dp$.

\subsection{Smooth Flow}

For $s \ne s_0$, the system obeys Hamilton's equations and the flow
is symplectic: $\mathcal{F}_t^*\omega = \omega$.

\subsection{Momentum Jump at the Fold}

The delta Lagrangian produces the impulse
\[
p(s_0^+) - p(s_0^-) \;=\; \mu\,n(s_0).
\]
Configuration $\gamma$ is continuous; momentum $p$ has a jump.

\subsection{The Fold as a Symplectic Map}

Define the fold map in phase space:
$\mathcal{F} : (\gamma, p) \mapsto (\gamma, p + \mu n)$.

\begin{theorem}[Symplectic Preservation]
$\mathcal{F}^*\omega = \omega$.
\end{theorem}

\begin{proof}
$d\gamma \wedge d(p + \mu n) = d\gamma \wedge dp
  + \mu\,d\gamma \wedge dn$.
Since $n$ depends only on $\gamma$, the term
$d\gamma \wedge dn = 0$.
Hence $\mathcal{F}^*\omega = \omega$.
\end{proof}

\subsection{Canonical Transformation}

The fold is generated by the distributional function
$S(\gamma) = \mu\,\Theta(|\kappa(\gamma)| - \kappa^*)$, so that
$p^+ = p^- + \partial S/\partial\gamma$.

\subsection{Full Transition}

Let $\Phi_t$ be the pre-fold Hamiltonian flow, $\mathcal{F}$ the fold
map, $\Psi_t$ the post-fold flow. The consistency condition from
Section~4 (basin structure) ensures the output of $\mathcal{F}$ lies
in the domain of $\Psi_t$. Then
\[
\mathcal{H} = \Psi_t \circ \mathcal{F} \circ \Phi_t
\]
is a piecewise-smooth symplectic map: $\mathcal{H}^*\omega = \omega$.

\subsection{Singularity Type and Momentum Structure}

\begin{itemize}
\item $A_1$ fold: single momentum jump.
\item $A_2$ cusp: momentum jump with first-order tangency constraint.
\item $A_3$ swallowtail: momentum jump with second-order tangency.
\end{itemize}

%-----------------------------------------------------------------------
\section{Connection to Generative Contact Mechanics}
\label{sec:connection}
%-----------------------------------------------------------------------

This section makes precise the relationship between the
singularity-theoretic framework developed in this volume and the
contact-geometric framework of \emph{Generative Contact Mechanics}
(Paper~1) and \emph{The dm3 Operator: Explicit Toy Model}
(Paper~2). We show that the operator sequence
$C \to K \to F \to U$
is the geometric precursor of the dm3 operator algebra, and that the
curvature threshold $\kappa^*$ corresponds canonically to the
embodiment threshold $\tau$.

\subsection{From Trajectories to Limit Cycles}

In this volume the central object is a trajectory
$\gamma : [0,T] \to (X,g)$ undergoing a localized generative
transition driven by curvature, injectivity loss, and stabilization.

In GCM the central object is a hyperbolic limit cycle
$\Gamma \subset S$ embedded in a dissipative dynamical system
satisfying Axioms~1--8 of the dm3 framework.

The conceptual shift is:
\begin{itemize}[leftmargin=*]
  \item \emph{Principia Orthogona}: analyzes single transitions
        along trajectories.
  \item \emph{GCM}: studies persistent post-transition attractors
        and their algebraic interactions.
\end{itemize}

The mathematical link is that every admissible generative transition
induces a locally attracting invariant set, and conversely, every
dm3 limit cycle admits a neighborhood whose formation is governed by
a fold-type transition.

\subsection{Curvature Threshold vs.\ Embodiment Threshold}

In this volume, folding is triggered when curvature reaches the
intrinsic geometric threshold
\[
  \kappa^*(x) = \frac{1}{\operatorname{foc}(x)},
\]
corrected by sectional curvature via the Rauch comparison theorem
(Section~\ref{sec:metric}).

In GCM, stochastic stability is governed by the embodiment threshold
\[
  \tau = \sqrt{c/\kappa},
\]
defined by the stochastic Lyapunov condition
$\mathcal{L}V \le -cV + \kappa\|\sigma\|^2$.

\begin{proposition}[Threshold Correspondence]\label{prop:threshold}
Let a generative transition produce a post-fold stabilized orbit
$\Gamma$ with Lyapunov function $V = d_g(\cdot,\Gamma)^2$. Then
the geometric condition $\kappa \uparrow \kappa^*$ is equivalent to
the dynamical condition that transverse contraction becomes dominant
over noise, yielding a finite embodiment threshold $\tau$.
\end{proposition}

\begin{proof}[Proof sketch]
As $\kappa \to \kappa^*$, the folding operator $F$ introduces
rank-1 loss and the post-fold unfolding $U$ selects a stable branch.
The resulting orbit $\Gamma$ has Floquet exponent
$\mu_{\max} < 0$ (transverse contraction). By the stochastic
Lyapunov condition of GCM, $\tau = \sqrt{c/\kappa}$ is finite
whenever $\mu_{\max} < 0$, i.e., precisely when curvature has
reached $\kappa^*$ and the fold has been completed. Thus
$\kappa^*$ is the geometric precursor of $\tau$: curvature
accumulation creates the conditions under which stochastic stability
becomes meaningful.
\end{proof}

\subsection{Fold Map vs.\ dm3 Operator}

The folding operator $F$ in this volume is a rank-1 singular map:
$\operatorname{rank}(dF) = \dim(X_C) - 1$, producing a finite
branch set classified by $A_1$--$A_3$ singularities.

In GCM, the dm3 operator
$\mathcal{A}_{\mathrm{dm3}} = \phi^{T^*/4}$
is the time-$(T^*/4)$ map of a smooth flow on the contact manifold
$M = S \times \mathbb{R}$, acting discretely along a limit cycle.

\begin{proposition}[Fold--dm3 Correspondence]\label{prop:fold}
The fold map $F$ is the piecewise-smooth, pre-contact limit of the
dm3 operator in the following sense.
\begin{enumerate}[label=\textup{(\roman*)}]
  \item $F$ enforces non-injectivity geometrically via rank loss;
        $\mathcal{A}_{\mathrm{dm3}}$ enforces recurrence dynamically
        via orbital stability.
  \item Under the contact extension $M = X \times \mathbb{R}$
        \textup{(}the construction of GCM, \S6\textup{)}, the
        impulsive momentum jump at the fold
        $p(s_0^+) - p(s_0^-) = \mu\,n(s_0)$
        corresponds to the contact Hamiltonian correction
        $H_{\mathrm{diss}} = -\gamma V e^{-\beta z}$,
        which drives trajectories toward $\Gamma$ off the orbit.
\end{enumerate}
\end{proposition}

\begin{proof}[Proof sketch]
Both the momentum jump and $H_{\mathrm{diss}}$ are corrections
that vanish on the invariant set and are nonzero off it. The
momentum jump is a distributional first-order correction in the
Hamiltonian sense; $H_{\mathrm{diss}}$ is its smooth
contact-geometric regularization as the contact manifold structure
is imposed via $M = X \times \mathbb{R}$.
\end{proof}

\subsection{Operator Grammar Correspondence}

The operator sequence of this volume maps onto the reduced dm3
operator grammar as follows.

\begin{center}
\begin{tabular}{ll}
\toprule
\emph{Principia Orthogona} & \emph{GCM / dm3} \\
\midrule
Compression $C$ & Basin contraction via $g_1$ \\
Curvature flow $K$ & Lyapunov descent $\dot{V} \le -cV$ \\
Fold $F$ & Contact bifurcation (Hopf, SN, NS) \\
Unfolding $U$ & Gradient flow to $\Gamma_{\mathrm{syn}}$ via $U_3$ \\
\bottomrule
\end{tabular}
\end{center}

The full dm3 grammar $g \to L \to R \to U$ extends this sequence
by adding multi-orbit coherence, resonance detection, and categorical
unification, which are intentionally absent from the present volume.

\subsection{Singularity Classification vs.\ Bifurcation Diagram}

This volume proves that admissible generative transitions are
precisely the Whitney singularities $A_1$, $A_2$, $A_3$.
Paper~2 proves that dm3 systems undergo exactly four bifurcations:
contact Hopf, saddle-node of limit cycles, Neimark--Sacker, and
slow-fast crossover.

\begin{remark}[Classification Correspondence]
The $A_1$--$A_3$ singularities classify the local geometry of dm3
bifurcations under the projection $M = S \times \mathbb{R} \to S$.
The fold $A_1$ corresponds to the contact Hopf and saddle-node
bifurcations (local rank loss in the radial direction); the cusp
$A_2$ corresponds to the Neimark--Sacker (loss of a second
direction at resonance); the swallowtail $A_3$ to the slow-fast
crossover (third-order degeneracy in the $z$-direction). Higher
singularities are excluded in both frameworks: here by the Morse
condition on $\Phi$; in dm3 by contact normal form rigidity
(Theorem~C of Paper~1).
\end{remark}

\subsection{Summary}

The three papers in this series have the following relationship.
\begin{itemize}[leftmargin=*]
  \item \textbf{This volume}: explains why folds must occur and
        classifies their local geometry.
  \item \textbf{Paper~1 (GCM)}: explains what stable structures
        exist after they occur, in the contact-geometric setting.
  \item \textbf{Paper~2 (dm3)}: proves that the resulting dynamics
        are computable, stable, and structurally closed.
\end{itemize}

No claim is made that every dm3 system arises from a literal
curvature fold in a physical manifold. The precise claim is:
every dm3 system admits a local geometric model whose generative
event is $\mathcal{A}$-equivalent to a fold-type transition governed
by a curvature threshold $\kappa^*$. This is supported by
Propositions~\ref{prop:threshold} and~\ref{prop:fold}.

%-----------------------------------------------------------------------
\section*{Closing Statement}
%-----------------------------------------------------------------------

This volume has developed a unified mathematical framework for
generative transitions. The central results are:
\begin{itemize}
\item a geometric foundation based on curvature, injectivity radius,
  and focal geometry;
\item a critical curvature threshold $\kappa^*$ defined intrinsically
  by the focal radius and corrected by the Rauch comparison theorem;
\item constructive definitions of operators $C$, $K$, $F$, $U$;
\item a free-discontinuity variational principle;
\item a Hamiltonian structure with a distributional generator at the
  fold;
\item a symplectic discontinuity preserving the canonical form;
\item a singularity classification restricted to the Whitney
  $A_1$--$A_3$ hierarchy; and
\item structural invariants holding across all admissible transitions.
\end{itemize}

The framework is placed within established mathematical traditions:
comparison geometry, singularity theory, geometric flows, variational
mechanics, and impulsive Hamiltonian systems.

This volume is deliberately limited in scope. It provides the
mathematical substrate on which cross-domain questions can be posed
rigorously. Volume Two constructs the contact-geometric realization
of this framework: the dm3 system, operator algebra, and four main
theorems including explicit structural stability with radius
$\varepsilon_0 = 1/3$. Volume Three instantiates the framework in
biology: allostatic stress, neural oscillations, circadian rhythms,
and immune adaptation cycles, with three falsifiable quantitative
predictions.

%-----------------------------------------------------------------------
\section*{Acknowledgments}
%-----------------------------------------------------------------------

The author acknowledges with gratitude the foundational influence of
the teachings of Paramahamsa Nithyananda and the yogic scriptural
tradition from which this work derives. The mathematical formalization
presented here is understood by the author as a translation of
principles encountered through that tradition.

This volume is the first in a series. The dm3 operator framework
developed in the companion papers \emph{Generative Contact Mechanics}
(Paper~1) and \emph{The dm3 Operator: Explicit Toy Model} (Paper~2)
constitutes the full mathematical realization of the generative
transition theory introduced here.
\cleardoublepage


%% ════════════════════════════════════════════════════════════
\part{Generative Contact Mechanics}
%% ════════════════════════════════════════════════════════════

\begin{center}
\emph{Dedicated to my children Vic~(R.I.P.), Giulia,
  Alice~(R.I.P.), Sarah~(R.I.P.), and David.}\\[4pt]
\emph{Once tiny, always strong.}
\end{center}

\vspace{1.5em}
\newpage

%-----------------------------------------------------------------------
\section*{Preface}
%-----------------------------------------------------------------------

This volume is the second in the Principia Orthogona series.
Volume~One~\cite{PO1} developed the singularity-theoretic and
variational foundations of generative transitions: the operator
sequence $C\to K\to F\to U$, the curvature threshold $\kappa^*$,
the Whitney $A_1$--$A_3$ singularity classification, and a
symplectic preservation theorem for the fold map.

The present volume constructs the explicit contact-geometric
realization of those foundations. The central results are:
\begin{itemize}[leftmargin=*]
  \item a precise correspondence between the geometric fold operator
        $F$ and the dm3 contact Hamiltonian dissipation
        (Section~\ref{sec:contact-fold}),
  \item a theorem establishing the equivalence of the curvature
        threshold $\kappa^*$ and the embodiment threshold $\tau$
        (Section~\ref{sec:threshold-equiv}),
  \item explicit verification of the correspondence in the dm3 toy
        model (Section~\ref{sec:verification}),
  \item a bifurcation analysis showing that the four dm3 bifurcations
        correspond to the singularity types of Volume~One
        (Section~\ref{sec:bifurcations}).
\end{itemize}

Throughout, we take as established the results of
Volume~One~\cite{PO1}, Generative Contact Mechanics~\cite{Paper1},
and the dm3 Toy Model~\cite{Paper2}. This volume does not re-derive
those results; it constructs the contact-geometric bridge between
the geometric framework of Volume~One and the dynamical framework
of~\cite{Paper1,Paper2}.

Volume~Three will instantiate the framework in biology:
allostatic stress, neural oscillations, circadian rhythms,
and immune adaptation cycles, with three falsifiable quantitative
predictions.

%-----------------------------------------------------------------------
\section{Introduction}
\label{sec:intro}
%-----------------------------------------------------------------------

\subsection{The problem: from geometric folds to dissipative dynamics}

Volume~One established that generative transitions are localized
geometric events classified by the Whitney $A_1$--$A_3$ hierarchy.
The fold operator $F$ was shown to act as a symplectic canonical
transformation on $T^*X$, and the full transition
$\cG = U\circ F\circ K\circ C$ was shown to be a piecewise-smooth
symplectic map. What Volume~One deliberately left open was the
question of \emph{post-fold stability}: once the fold has occurred
and the unfolding $U$ has selected a stable branch, what governs
the long-term dissipative dynamics near that branch?

The Hamiltonian framework of Volume~One cannot answer this question:
Liouville's theorem forbids attractors in symplectic systems on
compact manifolds. The answer is contact geometry. The contact
manifold $M = X\times\R$ with contact form $\alpha = dz - \lambda$
provides the correct geometric setting for dissipative dynamics
with limit cycle attractors, stochastic stability, and variational
structure. This is the framework of~\cite{Paper1}.

\subsection{Role of the dm3 toy model}
\label{subsec:toymodel-role}

The dm3 Toy Model~\cite{Paper2} is not an example among many,
nor a numerical illustration of general theory. Its role is
structural and methodological: it serves as a complete, explicit
realization of the abstract contact-geometric framework of~\cite{Paper1}
and as the bridge that certifies the realizability of the geometric
fold theory of Volume~One.

Its purpose can be stated precisely:

\begin{quote}
\emph{The dm3 Toy Model proves that the generative transition
framework is not merely definable, but fully instantiable by an
explicit, smooth, globally analyzable dynamical system.}
\end{quote}

This has three consequences. First, the definitions do not depend
on the toy model; no axiom is inferred from it and no theorem is
tailored to fit it. The model is introduced after the theory is
fully stated, to verify rather than to inspire. Second, the toy
model converts abstract structure into computable structure: every
operator, invariant, threshold, and bifurcation predicted by the
theory can be checked by direct calculation. Third, the toy model
certifies non-emptiness: without it, the axioms and operators could
be mutually consistent but never jointly satisfiable. The existence
of at least one concrete realization proves there are no hidden
contradictions.

The toy model does not claim to be generic, to represent a physical
system, or to capture all possible generative transitions. Its role
is logical, not empirical.

\subsection{Main results}

\begin{theorem}[A: Contact realization of the fold]
\label{thm:A}
The fold operator $F$ of Volume~One is the piecewise-smooth,
pre-contact limit of the dm3 operator
$\mathcal{A}_{\mathrm{dm3}} = \phi^{T^*/4}$ of~\cite{Paper1}.
Under the contact extension $M = X\times\R$, the impulsive
momentum jump $p^+ - p^- = \mu n$ at the fold corresponds to
the contact Hamiltonian correction
$H_{\mathrm{diss}} = -\gamma V e^{-\beta z}$
in the regularized limit $\beta\to\infty$.
\end{theorem}

\begin{theorem}[B: Threshold equivalence]
\label{thm:B}
Under the generative model of Volume~One and the dm3
realization of~\cite{Paper1}:
\[
  |\kappa|\uparrow\kappa^*
  \;\iff\;
  \mu_{\max} < 0
  \;\iff\;
  \tau = \sqrt{c/\kappa_{\mathrm{noise}}} \in (0,\infty).
\]
The curvature threshold $\kappa^*$ and the embodiment threshold
$\tau$ are two parameterizations of the same event: the onset of
transverse stability in the post-fold dissipative system.
\end{theorem}

\begin{theorem}[C: Singularity--bifurcation correspondence]
\label{thm:C}
The four bifurcations of the dm3 toy model~\cite{Paper2}
correspond to the Whitney singularity types of Volume~One
under the projection $M = S\times\R\to S$.
\end{theorem}

\subsection{Standing assumptions}

\begin{assumption}[Volume~One framework]
The operator sequence $C\to K\to F\to U$, the threshold $\kappa^*$,
and all results of~\cite{PO1} hold.
\end{assumption}

\begin{assumption}[dm3 framework]
The dm3 system, contact manifold $M=S\times\R$, and all results
of~\cite{Paper1,Paper2} hold.
\end{assumption}

%-----------------------------------------------------------------------
\section{Contact Realization of the Fold Operator}
\label{sec:contact-fold}
%-----------------------------------------------------------------------

\subsection{From Hamiltonian phase space to contact extension}

In Volume~One, the fold is realized in $T^*X$ as a symplectic
discontinuity. At fold point $s_0$:
\[
  p(s_0^+) - p(s_0^-) = \mu\,n(s_0),
\]
generated by $S(\gamma) = \mu\,\Theta(\kappa(\gamma)-\kappa^*)$
(\cite{PO1}, \S12). The fold map
$\mathcal{F}:(\gamma,p)\mapsto(\gamma,p+\mu n)$
preserves $\omega = d\gamma\wedge dp$
(\cite{PO1}, Theorem~11.1).

To capture post-fold stabilization, we pass to the contact
extension $M = X\times\R$ with $\alpha = dz - \lambda$,
$d\lambda=\omega$ (\cite{Paper1}, Definition~6.1).

\subsection{The fold as a contact discontinuity}

\begin{proposition}[Regularization of the fold generator]
\label{prop:regularization}
The contact dissipation Hamiltonian
$H_{\mathrm{diss}}(x,z) = -\gamma V(x)e^{-\beta z}$
is a smooth regularization of the distributional fold generator
$S(\gamma) = \mu\,\Theta(\kappa(\gamma)-\kappa^*)$
in the following precise sense: for each $\delta>0$, the
vector field correction $-\gamma\nabla V(x)e^{-\beta z}$
converges in the weak-$*$ topology on bounded measures to
$\mu\,\delta_{\{V=0\}}\otimes\delta_{\{z=0\}}$ as
$\beta\to\infty$, $\gamma\beta\to\mu/\norm{\nabla V}$.
\end{proposition}

\begin{proof}[Proof sketch]
Fix a test function $\phi\in C_c^\infty(X\times\R)$.
The integral $\int \gamma V(x)e^{-\beta z}\phi(x,z)\,dx\,dz$
concentrates near $\{V=0\}\times\{z=0\}$ as $\beta\to\infty$:
away from $z=0$ the factor $e^{-\beta z}\to 0$ uniformly on
$\{z\ge\delta\}$, and away from $\Orb=\{V=0\}$ the factor
$V(x)$ provides uniform decay. A Laplace-method argument shows
the limit equals $\mu\,\phi|_{\Orb\times\{0\}}$ with the
stated normalization $\gamma\beta\to\mu/\norm{\nabla V}$.
This is exactly the weak-$*$ action of the distributional
generator $S = \mu\,\Theta(\kappa-\kappa^*)$ at the fold
point, confirming that $H_{\mathrm{diss}}$ is the smooth
contact-geometric regularization of $S$.
\end{proof}

The contact dissipation has three structural properties:
(i)~$H_{\mathrm{diss}}|_\Orb=0$ (invariant set preserved);
(ii)~off $\Orb$: $H_{\mathrm{diss}}<0$ (contraction);
(iii)~$e^{-\beta z}$ weakens dissipation as action accumulates
(embodiment threshold).

\subsection{Relation to the dm3 normal form}

By \cite{Paper1} (Theorem~C), every dm3 system near $\Orb$ is
locally contact-diffeomorphic to
\begin{align}
  \dot\rho &= \mu_{\max}(1-e^{-\beta z})\rho + O(\rho^2),\\
  \dot\theta &= \omega + O(\rho),\\
  \dot z &= \omega - \abs{\mu_{\max}}\rho^2 e^{-\beta z} + O(\rho^3).
\end{align}
In the $C\to K\to F\to U$ sequence:
$\rho\to 0$ is the unfolding $U$;
$\mu_{\max}<0$ is the imprint of having crossed $\kappa^*$;
$z$ growing records accumulated dissipation from the fold.

\subsection{Correspondence table}

\begin{table}[ht]
\centering
\begin{tabular}{ll}
\toprule
\emph{Volume One (geometric)} & \emph{GCM / dm3 (contact)} \\
\midrule
Curvature threshold $\kappa^*$
  & Onset of transverse contraction \\
Fold impulse $p^+-p^-=\mu n$
  & Contact correction $H_{\mathrm{diss}}$ \\
Rank-1 Jacobian loss
  & Dissipative contact normal form \\
Unfolding $U$
  & Gradient flow to $\Orb$ via $U_3$~\cite{Paper1} \\
Distributional generator $S$
  & Regularized $H_{\mathrm{diss}}$
    (Proposition~\ref{prop:regularization}) \\
\bottomrule
\end{tabular}
\caption{Correspondence between Volume~One and the contact
  realization. This proves Theorem~A.}
\label{tab:correspondence}
\end{table}

%-----------------------------------------------------------------------
\section{Equivalence of \texorpdfstring{$\kappa^*$}{kappa*} and
  \texorpdfstring{$\tau$}{tau}}
\label{sec:threshold-equiv}
%-----------------------------------------------------------------------

\subsection{Geometric threshold implies stochastic threshold}

\begin{lemma}[Fold activation produces hyperbolicity]
\label{lem:fold-hyp}
If $K$ drives curvature to $\kappa^*$, $F$ is rank-1, and $U$
selects a nondegenerate branch, then $\Orb$ is hyperbolic with
$\mu_{\max}<0$ and $\dot V\le -cV$ for some $c>0$.
\end{lemma}

\begin{proof}[Proof sketch]
Rank-1 loss at $F$ and Morse nondegeneracy of $\Phi$ yield
transverse contraction. Floquet theory gives $\mu_{\max}<0$
(\cite{PO1}, Theorem~3.3; \cite{Paper1}, Proposition~3.2).
\end{proof}

\begin{theorem}[Geometric implies stochastic threshold]
\label{thm:geom-stoch}
Under Lemma~\ref{lem:fold-hyp},
$\cL V\le -cV+\kappa_{\mathrm{noise}}\norm{\sigma}^2$
and $\tau=\sqrt{c/\kappa_{\mathrm{noise}}}\in(0,\infty)$.
\end{theorem}

\begin{proof}
$\dot V\le -cV$ plus the It\^{o} correction
$\frac{1}{2}\norm{\sigma}^2\norm{\Hess V}$ gives
$\cL V\le -cV+\kappa_{\mathrm{noise}}\norm{\sigma}^2$
with $\kappa_{\mathrm{noise}}=\frac{1}{2}\sup\norm{\Hess V}$.
Then $\tau=\sqrt{c/\kappa_{\mathrm{noise}}}>0$
(\cite{Paper1}, Definition~3.4).
\end{proof}

\subsection{Stochastic threshold implies geometric threshold}

\begin{lemma}[Finite $\tau$ implies transverse contraction]
\label{lem:tau-contract}
If $\cL V\le -cV+\kappa_{\mathrm{noise}}\norm{\sigma}^2$ with
$c>0$, then $\mu_{\max}<0$ and $\dot V\le -cV$ near $\Orb$.
\end{lemma}

\begin{proof}
Exponential decay of $\mathbb{E}[V(X_t)]$ forces $\mu_{\max}<0$
by stochastic stability theory
(\cite{Paper1}, Propositions~3.2 and~3.4).
\end{proof}

\begin{theorem}[Stochastic implies geometric threshold]
\label{thm:stoch-geom}
If $\tau\in(0,\infty)$, then the trajectory must have crossed
$\kappa^*$ and undergone a rank-1 fold.
\end{theorem}

\begin{proof}
By Lemma~\ref{lem:tau-contract}, $\mu_{\max}<0$, so a hyperbolic
attracting cycle exists. Suppose $|\kappa|<\kappa^*$ everywhere.
Then \cite{PO1} (Invariant~7.5) gives injectivity before threshold:
no rank loss, no mechanism for a hyperbolic attracting cycle.
Contradiction. So $\kappa=\kappa^*$ must be reached.
\end{proof}

\subsection{The equivalence theorem}

\begin{theorem}[Threshold equivalence]
\label{thm:equiv}
\[
  |\kappa|\uparrow\kappa^*
  \;\iff\;
  \mu_{\max}<0
  \;\iff\;
  \tau\in(0,\infty).
\]
$\kappa^*$ is the geometric precursor of $\tau$: curvature
accumulation creates the conditions under which stochastic
stability becomes meaningful. This proves Theorem~B.
\end{theorem}

\begin{proof}
Forward: Theorem~\ref{thm:geom-stoch}.
Backward: Theorem~\ref{thm:stoch-geom}.
Middle: $\tau>0\iff c>0\iff\mu_{\max}<0$
(Lemma~\ref{lem:tau-contract}).
\end{proof}

\begin{remark}[Scope]
The equivalence is local to the fold neighborhood and the
post-fold tubular neighborhood of $\Orb$. The exclusion of
$A_k$ ($k\ge 4$) in Volume~One (Morse condition) aligns with
contact normal form rigidity in~\cite{Paper1} (Theorem~C).
\end{remark}

%-----------------------------------------------------------------------
\section{Explicit Verification in the dm3 Toy Model}
\label{sec:verification}
%-----------------------------------------------------------------------

\subsection{The system}

\subsection{The verification theorem}
\label{subsec:verification-theorem}

Before exhibiting the explicit computations, we state the logical
content of the entire verification as a single theorem.

\begin{theorem}[Verification of GCM by the dm3 Toy Model]
\label{thm:verification}
There exists an explicit smooth dynamical system on a contact
manifold---the dm3 toy model of~\cite{Paper2}---in which all of
the following are verified by direct computation:
\begin{enumerate}[label=\textup{(\arabic*)}]
  \item \textbf{Axiom consistency.} All eight dm3 axioms
        of~\cite{Paper1} are satisfied simultaneously.
        The axiom class is non-empty.
  \item \textbf{Contact structure.} The system lives on an explicit
        contact manifold with $\alpha=dz-r^2\,d\theta$; all contact
        identities hold by direct verification.
  \item \textbf{Contact normal form.} The system is explicitly
        contact-diffeomorphic to the universal normal form of
        Theorem~C of~\cite{Paper1} with $(\mu_{\max},\omega,\beta)=(-2,1,1)$.
        The normal form is non-empty.
  \item \textbf{Operator algebra closure.} Every operator ($g$-,
        $L$-, $R$-, $U$-, $B$-operators) is instantiated explicitly
        and preserves the dm3 axioms. The full grammar closes.
  \item \textbf{Stochastic stability.} $\tau=2$ is computed in
        closed form. The stationary Fokker--Planck measure is
        Gaussian, concentrating on $\Orb$ for $\sigma<\tau$ and
        spreading for $\sigma>\tau$.
  \item \textbf{Global dynamics.} The global attractor is $\Orb_{12}$;
        all invariant sets are identified; all four predicted
        bifurcations are exhibited.
\end{enumerate}
\end{theorem}

\begin{proof}
Items (1)--(6) are verified in~\cite{Paper2} by direct computation:
(1)~\S2.4; (2)~Prop.~2.2; (3)~\S7; (4)~\S3; (5)~\S5 and Thm.~D;
(6)~\S\S4,6 and Thm.~A.
\end{proof}

\begin{remark}[Logical status of the verification]
The dm3 Toy Model $\mathcal{D}_0$ provides a constructive
verification of GCM in the strongest possible sense: not by
example or illustration, but by explicit realization. All abstract
components of GCM---axioms, geometry, operators, thresholds,
normal forms, stochastic predictions, and global dynamics---are
shown to coexist within a single smooth contact-geometric system.

In one sentence: \emph{the dm3 Toy Model is a complete verification
witness for Generative Contact Mechanics, demonstrating that the
entire abstract framework is jointly realizable, computable, and
dynamically consistent.}

The toy model does not claim to be generic or to represent a
physical system. Its role is logical, not empirical.
\end{remark}

\subsection{The system}
\label{subsec:system}

On $M=\R^2_{>0}\times\R$ with polar coordinates $(r,\theta,z)$:
\begin{align}
  \dot r &= r(1-r^2)+2(r-1)e^{-z},\label{eq:rdot}\\
  \dot\theta &= 1,\label{eq:tdot}\\
  \dot z &= r^2 - 2(r-1)^2 e^{-z}.\label{eq:zdot}
\end{align}
Limit cycle: $\Orb=\{r=1\}$, $T^*=2\pi$.
Contact form: $\alpha=dz-r^2\,d\theta$
(\cite{Paper2}, Proposition~2.2).

\subsection{Threshold verification}

\begin{proposition}[Explicit threshold values]
\label{prop:explicit}
$\mu_{\max}=-2$, $\kappa_{\mathrm{noise}}=1$, $\tau=2$.
The transverse eigenvalue $\lambda(z)=-2(1-e^{-z})$ satisfies:
$\lambda(0)=0$ (neutral, pre-embodiment),
$\lambda(z)<0$ for $z>0$ (attracting, post-embodiment),
$\lambda(z)\to -2$ as $z\to\infty$ (full dm3 rate).
\end{proposition}

\begin{proof}
Linearize at $r=1$: $\dot\rho=-2(1-e^{-z})\rho+O(\rho^2)$.
Generator: $\cL V=-4V(1-e^{-z})+\sigma^2$, giving $c\to 4$,
$\kappa_{\mathrm{noise}}=1$, $\tau=2$
(\cite{Paper2}, Lemma~2.1 and Proposition~2.4).
\end{proof}

The neutral stability at $z=0$ is the mathematical content of
the embodiment threshold: the orbit earns its stability by
accumulating action. The crossing $z=0\to z>0$ in the contact
variable corresponds exactly to the curvature crossing
$\kappa\to\kappa^*$ in Volume~One.

\subsection{Normal form verification}

\begin{proposition}[Contact normal form]
\label{prop:nf}
In coordinates $(\rho,\theta,z)$ with $\rho=r-1$, the system
takes the form
\[
  \dot\rho=-2(1-e^{-z})\rho+O(\rho^2),\quad
  \dot\theta=1+O(\rho),\quad
  \dot z=1-2\rho^2 e^{-z}+O(\rho^3),
\]
the contact normal form of~\cite{Paper1} (Theorem~C) with
$(\mu_{\max},\omega,\beta)=(-2,1,1)$.
\end{proposition}

\subsection{Stability radius}

The structural stability radius of Theorem~D of~\cite{Paper1} is
\[
  \varepsilon_0
  = \frac{\abs{\mu_{\max}}}{2(1+\sup_\Orb\norm{\Hess V}_\infty)},
\]
where $\abs{\mu_{\max}}$ denotes the magnitude of the maximal
transverse Floquet exponent --- a spectral quantity --- not the
Lyapunov rate $c$ appearing in the stochastic Lyapunov condition
$\cL V \le -cV + \kappa_{\mathrm{noise}}\norm{\sigma}^2$.

\begin{remark}[Convention: Floquet exponent, not Lyapunov rate]
\label{rem:convention}
In the toy model $V=(r-1)^2$, the Lyapunov rate satisfies
$\dot V = -4V$ near $\Orb$, giving $c=4$, while the Floquet
exponent is $\mu_{\max}=-2$. These differ by a factor of $2$
arising from the quadratic structure of $V$. The formula for
$\varepsilon_0$ in Theorem~D of~\cite{Paper1} uses
$\abs{\mu_{\max}}=2$ (the spectral quantity), not $c=4$.
This gives
\[
  \varepsilon_0
  = \frac{2}{2(1+2)}
  = \frac{1}{3}.
\]
This convention is fixed throughout the series. Using $c$ in
place of $\abs{\mu_{\max}}$ would give $\varepsilon_0=2/3$ and
would be a category error: the stability radius is determined by
the spectrum of the linearized flow, not by the rate of Lyapunov
descent.
\end{remark}

%-----------------------------------------------------------------------
\section{Singularity--Bifurcation Correspondence}
\label{sec:bifurcations}
%-----------------------------------------------------------------------

\subsection{The four bifurcations of the toy model}

From \cite{Paper2} (Theorem~C):
\begin{enumerate}[label=(\roman*)]
  \item \textbf{Contact Hopf} at $\gamma=e^{z_0}$: limit cycle
        loses stability, new cycle bifurcates.
  \item \textbf{Saddle-node} at $\eta\approx 0.15$: two cycles
        collide, multi-orbit structure collapses.
  \item \textbf{Neimark--Sacker} at detuning $|\Delta|=\Delta^*$:
        resonant orbit loses stability, 2-torus bifurcates.
  \item \textbf{Slow-fast crossover} at $\beta=\beta^*$: smooth
        transition between contact and classical regimes.
\end{enumerate}

\subsection{Correspondence with Whitney singularities}

\begin{proposition}[Singularity--bifurcation correspondence]
\label{prop:sing-bif}
Under projection $M=S\times\R\to S$:
\begin{center}
\begin{tabular}{lll}
\toprule
Singularity & Codimension & dm3 bifurcation \\
\midrule
$A_1$ (fold) & 0 & Contact Hopf and saddle-node \\
$A_2$ (cusp) & 1 & Neimark--Sacker \\
$A_3$ (swallowtail) & 2 & Slow-fast crossover \\
\bottomrule
\end{tabular}
\end{center}
Higher singularities excluded: in Volume~One by the Morse
condition; in~\cite{Paper1} by contact normal form rigidity
(Theorem~C).
\end{proposition}

\begin{proof}[Proof sketch]
$A_1$: rank-1 loss in one transverse direction --- mechanism of
both Hopf (radial stability loss) and saddle-node (radial
collision). $A_2$: rank-1 loss with second-order tangency ---
matches Neimark--Sacker (second angular direction destabilizes
at resonance). $A_3$: rank-1 with third-order tangency ---
matches slow-fast crossover (third-order change in $z$-direction).
\end{proof}

This proves Theorem~C.

%-----------------------------------------------------------------------
\section{Discussion}
\label{sec:discussion}
%-----------------------------------------------------------------------

\subsection{Summary of the three main results}

\begin{enumerate}[label=(\arabic*)]
  \item The fold operator $F$ is the geometric precursor of the
        dm3 contact Hamiltonian $H_{\mathrm{diss}}$; the
        distributional generator $S$ is regularized by
        $H_{\mathrm{diss}}$ in the limiting sense of
        Proposition~\ref{prop:regularization} (Theorem~A).
  \item $\kappa^*$ and $\tau$ are equivalent: each is finite
        and positive if and only if the other is, both detecting
        the onset of transverse stability (Theorem~B).
  \item The four dm3 bifurcations correspond to Whitney
        $A_1$--$A_3$ types, completing the singularity-theoretic
        classification of the toy model dynamics (Theorem~C).
\end{enumerate}

\subsection{Role in the trilogy}

\begin{itemize}[leftmargin=*]
  \item \textbf{Volume~One}: why folds must occur;
        local geometric classification.
  \item \textbf{Volume~Two}: what stable structures exist after
        folds; contact-geometric realization.
  \item \textbf{Volume~Three}: domain-specific instantiations
        (plasma, markets, neural geometry).
\end{itemize}

\subsection{Open problems}

\begin{enumerate}[label=(\arabic*)]
  \item \textbf{Global equivalence.} Theorem~B is local.
        A global version requires showing every $\tau$-stable
        dm3 system arises from a fold transition globally on $X$.
  \item \textbf{Higher resonances.} Systematic treatment of
        $k$:$m$ correspondence between higher $A_k$ singularities
        and higher resonances.
  \item \textbf{Volume~Three.} Domain-specific $(X,g,\Phi)$
        with explicit computation of $\kappa^*$ and $\tau$ from
        data.
\end{enumerate}

%-----------------------------------------------------------------------
\section*{Acknowledgments}
%-----------------------------------------------------------------------

The author acknowledges with gratitude the foundational influence
of the teachings of Paramahamsa Nithyananda and the yogic
scriptural tradition from which this work derives. The mathematical
formalization presented here is understood by the author as a
translation of principles encountered through that tradition.

%-----------------------------------------------------------------------
\begin{thebibliography}{99}
%-----------------------------------------------------------------------

\bibitem{PO1}
P.~Nogueira~Grossi,
\emph{Principia Orthogona, Volume One: The Mathematics of
  Generative Transitions},
preprint, HAL, 2026.

\bibitem{Paper1}
P.~Nogueira~Grossi,
\emph{Generative Contact Mechanics: A Geometric Framework for
  Dissipative Systems with Structured Limit Cycles},
submitted to J.\ Geom.\ Mech., 2026.

\bibitem{Paper2}
P.~Nogueira~Grossi,
\emph{The dm3 Operator: Explicit Toy Model and Global Dynamical
  Analysis},
submitted to SIAM J.\ Appl.\ Dyn.\ Syst., 2026.

\bibitem{Arnold}
V.~I.~Arnold,
\emph{Mathematical Methods of Classical Mechanics},
2nd ed., Springer, New York, 1989.

\bibitem{Bravetti2017}
A.~Bravetti,
\emph{Contact Hamiltonian mechanics},
Ann.\ Phys.\ \textbf{376} (2017), 17--39.

\bibitem{deLeonLainz2019}
M.~de~Le{\'o}n and M.~Lainz~Valc{\'a}zar,
\emph{Contact Hamiltonian systems},
J.\ Math.\ Phys.\ \textbf{60} (2019), 102902.

\bibitem{GH}
J.~Guckenheimer and P.~Holmes,
\emph{Nonlinear Oscillations, Dynamical Systems, and Bifurcations
  of Vector Fields},
Springer, New York, 1983.

\bibitem{HPS}
M.~W.~Hirsch, C.~C.~Pugh, and M.~Shub,
\emph{Invariant Manifolds},
Lecture Notes in Mathematics 583, Springer, Berlin, 1977.

\bibitem{Hasminski}
R.~Z.~Has{'}mi{\u{n}}ski{\u{\i}},
\emph{Stochastic Stability of Differential Equations},
Sijthoff \& Noordhoff, 1980.

\bibitem{Shub}
M.~Shub,
\emph{Global Stability of Dynamical Systems},
Springer, New York, 1987.

\end{thebibliography}
\cleardoublepage


%% ════════════════════════════════════════════════════════════
\part{The dm3 Operator: Explicit Toy Model}
%% ════════════════════════════════════════════════════════════

\title[Generative Contact Mechanics]{%
  Generative Contact Mechanics:\\
  A Geometric Framework for Dissipative Systems\\
  with Structured Limit Cycles}

\author{Pablo Nogueira Grossi}
\address{Academic Coordinator, UCEDA School (ESL), Elizabeth, NJ\\
Newark, NJ, USA}
\email{pablogrossi@hotmail.com}

\keywords{contact geometry, limit cycles, dissipative dynamical systems,
  operator algebra, structural stability, resonance, stochastic stability,
  nonequilibrium dynamics, Lyapunov functions}

\subjclass[2020]{37C10, 37C27, 37D30, 53D10, 60H10, 18A30}



}

%%=============================================================
\section{Introduction}\label{sec:intro}
%%=============================================================

\subsection{The problem: dissipative systems with structured limit cycles}

Dissipative dynamical systems with hyperbolic limit cycles appear throughout
the natural sciences---in biology as oscillatory cellular and neural
processes, in physics as driven nonequilibrium systems, in engineering as
controlled oscillators, and in mathematics as paradigmatic examples of
nonlinear dynamics. The classical theory provides excellent tools for
analyzing a single limit cycle in isolation: Floquet theory characterizes
transverse stability, Lyapunov theory gives quantitative convergence rates,
and stochastic extensions handle noise.

What is less well understood is how limit cycles interact, unify, and
organize into multi-scale coherent structures while remaining robust under
perturbation. This paper develops a geometric framework---generative contact
mechanics---for precisely this class of behavior. The central insight is
that the correct geometric home for dissipative systems with structured limit
cycles is not symplectic geometry (which forbids attractors) but contact
geometry, where dissipation is encoded in the geometry itself.

\subsection{The contact geometric approach}

A hyperbolic limit cycle is an attractor: nearby trajectories converge to it
exponentially. This is incompatible with Hamiltonian mechanics, where
Liouville's theorem guarantees that the phase-space volume is preserved and
no attractors can exist on compact symplectic manifolds. Standard Lagrangian
mechanics is similarly insufficient: the Euler--Lagrange equations are
time-reversible and do not admit attractors without additional dissipation
terms added by hand.

Contact geometry provides a natural resolution. A contact manifold
$(M^{2n+1},\alpha)$ carries a Reeb vector field and a contact Hamiltonian
formalism in which dissipation is built into the geometric structure: the
contact Hamiltonian $H$ evolves along trajectories rather than remaining
constant, and the action variable $z\in\R$ encodes accumulated dissipation.
This makes contact geometry the natural setting for dissipative systems
with limit cycle attractors, stochastic stability thresholds, and
variational structure~\cite{Bravetti2017,deLeonLainz2019}.

The present paper formalizes this setting through the dm3 system and
operator algebra, and shows that the resulting framework is: (a)~geometrically
natural (contact normal form, Theorem~C), (b)~categorically coherent
(closure of $\dmcat$ under unification, Theorem~B), and (c)~physically
robust ($C^1$-structural stability, Theorem~D).

\subsection{Main results}

The four main theorems are stated informally here and proved in later
sections. Precise statements appear at the end of each relevant section.

\begin{theorem}[A: dm3 existence and stability]\label{thm:A}
Let $\dm=(S,g,Y,\Orb,V,\sigma)$ satisfy
Axioms~\ref{ax:orbit}--\ref{ax:hyperbolic}. Then:
\begin{enumerate}[label=\textup{(\roman*)}]
  \item $\Orb$ is a hyperbolic limit cycle of $Y$ with minimal period $T^*>0$.
  \item $\Orb$ is exponentially orbitally stable with rate $c=\abs{\mu_{\max}}$.
  \item For noise amplitude below the embodiment threshold $\tau$,
        $\mathbb{E}[V(X_t)]\to 0$.
  \item The winding integral $\oint_{\Orb}\lambda$ is a topological invariant
        preserved under all admissible deformations.
        \textup{(Proved in Proposition~\ref{prop:winding}.)}
\end{enumerate}
\end{theorem}

\begin{theorem}[B: closure under unification]\label{thm:B}
Let $\dm_1,\dm_2\in\dmcat$ be in $k$:$m$ resonance with admissible span.
Then:
\begin{enumerate}[label=\textup{(\roman*)}]
  \item The unification operator $\cU=U_3\circ U_2\circ U_1\circ R_3\circ R_2\circ R_1$
        is well-defined and $\cU(\dm_1,\dm_2)\in\dmcat$.
  \item The canonical invariants satisfy
        $T^*_{12}=kT^*_1/\gcd(k,m)$,\;
        $\mu_{\max,12}=\max(\mu_{\max,i})$,\;
        $\tau_{12}\le\min(\tau_1,\tau_2)$.
  \item \textbf{Unification weakens noise tolerance:}
        $\tau_{12}\le\min(\tau_1,\tau_2)$.
\end{enumerate}
\end{theorem}

\begin{theorem}[C: contact normal form]\label{thm:C}
Let $\dm$ satisfy Axioms~\ref{ax:orbit}--\ref{ax:hyperbolic}. In a tubular
neighborhood of $\Orb$ in $(M,\alpha)$, there exist contact coordinates
$(\rho,\theta,z)$ with $\Orb=\{\rho=0\}$ such that
\begin{align*}
  \dot\rho &= \mu_{\max}(1-e^{-\beta z})\rho + O(\rho^2),\\
  \dot\theta &= \omega + O(\rho),\\
  \dot z &= \omega - \abs{\mu_{\max}}\rho^2 e^{-\beta z} + O(\rho^3).
\end{align*}
Every dm3 system is locally contact-diffeomorphic to this normal form.
The parameters $(\mu_{\max},\omega,\beta)$ are the canonical invariant
triple.
\end{theorem}

\begin{theorem}[D: structural stability]\label{thm:D}
Let $\dm$ satisfy Axioms~\ref{ax:orbit}--\ref{ax:hyperbolic} with
$\mu_{\max}<0$. Define
\[
  \varepsilon_0 \coloneqq
  \frac{\abs{\mu_{\max}}}{2\bigl(1+\sup_{\Orb}\norm{\Hess V}_\infty\bigr)}.
\]
For any $C^1$ perturbation $F$ with $\norm{F}_{C^1}<\varepsilon_0$, and
any finite admissible operator sequence within parameter bounds, the output
$\dm'$ satisfies Axioms~\ref{ax:orbit}--\ref{ax:hyperbolic} with: a unique
hyperbolic limit cycle $\Orb'$ near $\Orb$; the same winding number;
canonical invariants varying by $O(\varepsilon_0)$; and boundaries $B_i'$
varying $C^1$-smoothly.
\end{theorem}

\subsection{Relation to existing work}\label{sec:related}

\emph{Contact Hamiltonian mechanics.}
Bravetti~\cite{Bravetti2017} and de~Le\'on--Lainz~\cite{deLeonLainz2019}
developed the modern formalism of contact Hamiltonian mechanics, showing
that contact geometry provides a variational framework for dissipative
systems. Gaset et al.~\cite{Gaset2020} extended this to Lagrangian
contact mechanics. The present paper differs from these works in three
respects: (a)~we focus on limit cycle attractors rather than fixed-point
attractors; (b)~we construct a categorical operator algebra acting on
contact systems; and (c)~we prove structural stability of the operator
algebra with an explicit stability radius.

\emph{Metriplectic systems.}
Morrison~\cite{Morrison1986} introduced the metriplectic framework for
systems combining Hamiltonian and dissipative dynamics. The dm3 contact
formulation is related but distinct: metriplectic structure separates the
Poisson and metric brackets, while our contact structure encodes dissipation
in the contact form itself through the action variable $z$.

\emph{Normally hyperbolic invariant manifolds.}
Hirsch, Pugh, and Shub~\cite{HPS} established the persistence and
$C^r$-regularity of normally hyperbolic invariant manifolds under
perturbation. Theorem~D uses this theory as a foundational input.
The novelty is combining NHIM persistence with the operator algebra
and establishing the quantitative stability radius $\varepsilon_0$.

\emph{Arnold tongues and coupled oscillators.}
The resonance structure of the $R$-operators (period resonance, Arnold
tongues, mode locking) is classical~\cite{Arnold,Pikovsky}. The contribution
here is embedding these phenomena in a categorical structure: resonance is
the precondition for a well-defined pushout in $\dmcat$. The prediction
$\tau_{12}\le\min(\tau_i)$ is new.

\subsection{Organization}

Section~\ref{sec:background} fixes notation and background.
Section~\ref{sec:dm3} defines dm3 systems and proves Theorem~A.
Section~\ref{sec:operators} constructs the operator algebra and proves
Theorem~B.
Section~\ref{sec:boundary} establishes boundary operators and the
operator grammar.
Section~\ref{sec:contact} develops the contact physics and proves
Theorem~C.
Section~\ref{sec:stability} proves Theorem~D.
Section~\ref{sec:discussion} discusses open problems and connections.
Appendices give the categorical and lemma details.

%%=============================================================
\section{Background}\label{sec:background}
%%=============================================================

\subsection{Riemannian manifolds and flows}

Throughout, $(S,g)$ denotes a smooth, connected, complete Riemannian
manifold of dimension $n\ge 2$. We write $T_xS$ for the tangent space
at $x$, $\mathfrak{X}(S)$ for smooth vector fields, and $\Omega^k(S)$
for smooth $k$-forms. For $X\in\mathfrak{X}(S)$ complete, the flow
$\phi^t\colon S\to S$ satisfies $\phi^0=\mathrm{id}$ and
$\tfrac{d}{dt}\phi^t(x)=X(\phi^t(x))$.

\subsection{Hyperbolic limit cycles and Floquet theory}

\begin{definition}\label{def:limitcycle}
A \emph{limit cycle} of $Y\in\mathfrak{X}(S)$ is a compact connected
invariant submanifold $\Orb\cong S^1$ with minimal period $T^*>0$.
It is \emph{hyperbolic} if all Floquet multipliers except the trivial
multiplier $\mu_0=1$ satisfy $\abs{\mu_j}<1$ for $j=1,\ldots,n-1$.
The \emph{maximal transverse Lyapunov exponent} is
$\mu_{\max}\coloneqq\max_j\tfrac{1}{T^*}\log\abs{\mu_j}<0$.
\end{definition}

The stable manifold theorem~\cite{HPS} applied to $\Orb$ gives: if
$\Orb$ is hyperbolic, then $\partial\cB(\Orb)=\mathcal{W}^s(\Lambda)$
for some invariant set $\Lambda\subset S\setminus\cB(\Orb)$.

\subsection{Lyapunov functions}

\begin{definition}\label{def:lyapunov}
A smooth function $V\colon S\to\R_{\ge 0}$ is a \emph{Lyapunov function}
for $\Orb$ if $V(x)=0\iff x\in\Orb$ and
$\dot V(x)\coloneqq g(\nabla V(x),Y(x))<0$ for $x\notin\Orb$.
The standard choice is $V=\dg(\cdot,\Orb)^2$.
\end{definition}

\subsection{Stochastic differential equations}

For the It\^o SDE $dX_t=Y(X_t)\,dt+\sigma(X_t)\,dW_t$ on $S$,
the generator is
\[
  \cL f(x) = g(\nabla f(x),Y(x))
             + \tfrac12\mathrm{tr}\bigl(\sigma(x)^\top\Hess_f(x)\,\sigma(x)\bigr).
\]
Standard stochastic Lyapunov theory~\cite{Hasminski} gives: if
$\cL V\le -cV+\kappa\norm\sigma^2$ for $c,\kappa>0$, then
$\mathbb{E}[V(X_t)]\to 0$ for noise amplitude below $\tau=\sqrt{c/\kappa}$.

\subsection{Contact manifolds}

\begin{definition}\label{def:contact}
A \emph{contact manifold} $(M^{2n+1},\alpha)$ is a smooth manifold
with a 1-form $\alpha$ satisfying $\alpha\wedge(d\alpha)^n\ne 0$
everywhere. The \emph{Reeb vector field} $R_\alpha$ is the unique
vector field with $\iota_{R_\alpha}d\alpha=0$ and $\alpha(R_\alpha)=1$.
For a smooth $H\colon M\to\R$, the \emph{contact Hamiltonian vector
field} $X_H$ is the unique vector field satisfying
$\iota_{X_H}d\alpha=dH-(R_\alpha H)\alpha$ and $\alpha(X_H)=-H$.
\end{definition}

See Geiges~\cite{Geiges} and Bravetti~\cite{Bravetti2017} for background.

%%=============================================================
\section{The dm3 operator}\label{sec:dm3}
%%=============================================================

\subsection{Definition}

\begin{definition}[dm3 system]\label{def:dm3}
A \emph{dm3 system} is a tuple $\dm=(S,g,Y,\Orb,V,\sigma)$ where
$(S,g)$ is a smooth complete Riemannian manifold, $Y\in\mathfrak{X}(S)$
is a smooth vector field with complete flow $\phi^t$, $\Orb\subset S$
is a hyperbolic limit cycle, $V\colon S\to\R_{\ge 0}$ is a smooth
Lyapunov function for $\Orb$, and $\sigma\colon S\to\R^{n\times m}$
is a smooth noise coefficient. The \emph{dm3 operator} is
$\mathcal{A}_{\mathrm{dm3}}\coloneqq\phi^T$ for $T=T^*/4$.
\end{definition}

\subsection{The eight axioms}\label{subsec:axioms}

\begin{axiom}[Orbit identity]\label{ax:orbit}
There exist four distinct points $p_I,p_{VI},p_{III},p_{VII}\in\Orb$
traversed cyclically:
$\phi^{T^*/4}(p_I)=p_{VI}$, $\phi^{T^*/4}(p_{VI})=p_{III}$,
$\phi^{T^*/4}(p_{III})=p_{VII}$, $\phi^{T^*/4}(p_{VII})=p_I$.
\end{axiom}

\begin{axiom}[Generative closure]\label{ax:closure}
$\Orb$ is compact, connected, and invariant:
$\phi^t(\Orb)=\Orb$ for all $t\in\R$.
\end{axiom}

\begin{axiom}[Structural primacy]\label{ax:primacy}
There exists a neighborhood $\cN(\Orb)$ such that
$\dot V(x)=g(\nabla V(x),Y(x))<0$ for all $x\in\cN(\Orb)\setminus\Orb$.
\end{axiom}

\begin{axiom}[Perturbation response]\label{ax:response}
There exist $c>0$ and $\cN(\Orb)$ such that
$\dot V(x)\le -c\,V(x)$ for all $x\in\cN(\Orb)$.
\end{axiom}

\begin{axiom}[Noise-as-fuel]\label{ax:noise}
The SDE generator satisfies
$\cL V(x)\le -c\,V(x)+\kappa\norm{\sigma(x)}^2$
for constants $c,\kappa>0$.
\end{axiom}

\begin{axiom}[Non-collapse]\label{ax:noncollapse}
$Y(x)\ne 0$ for all $x\in\Orb$, and $\phi^t$ has minimal period
$T^*>0$ on $\Orb$.
\end{axiom}

\begin{axiom}[Non-coercion]\label{ax:smooth}
$Y$ is $C^\infty$ on $S$. No discontinuous projection, clamping,
or hard boundary enforcement appears in the definition of $Y$.
\end{axiom}

\begin{axiom}[Hyperbolicity]\label{ax:hyperbolic}
All Floquet multipliers of $\Orb$ except the trivial one satisfy
$\abs{\mu_j}<1$ for $j=1,\ldots,n-1$.
\end{axiom}

\begin{remark}
These eight axioms are logically independent. The original ten axioms
are reduced here by: merging the stochastic pair into
Axiom~\ref{ax:noise}; and moving invariance under the $g$- and
$L$-operators and multi-scale compatibility to their respective
sections as theorems.
\end{remark}

\subsection{Orbital stability}

\begin{proposition}[Exponential orbital stability]\label{prop:orbital}
Let $\dm$ satisfy Axioms~\ref{ax:orbit}--\ref{ax:hyperbolic}. Then
there exist $c>0$ and $\cN(\Orb)$ such that
$\dot V(x)\le -c\,V(x)$ for all $x\in\cN(\Orb)$.
\end{proposition}

\begin{proof}
Axiom~\ref{ax:hyperbolic} gives $\mu_{\max}<0$. By Floquet theory
(see~\cite{GH,Shub}), the linearized transverse flow contracts with
rate $\abs{\mu_{\max}}$. With $V=\dg(\cdot,\Orb)^2$, we get
$\dot V\le 2\mu_{\max}V+O(V^{3/2})$ in a tubular neighborhood.
Taking $\cN(\Orb)$ small enough absorbs the higher-order term,
giving $\dot V\le -c\,V$ with $c=\abs{\mu_{\max}}$.
\end{proof}

\subsection{Embodiment threshold and stochastic stability}

\begin{proposition}[Stochastic Lyapunov condition]\label{prop:stoch}
If Axiom~\ref{ax:noise} holds with constants $c,\kappa$, then for
noise amplitude $\norm\sigma<\tau\coloneqq\sqrt{c/\kappa}$:
\[
  \mathbb{E}[V(X_t)]\le e^{-ct}V(X_0)+\frac{\kappa}{c}\norm\sigma^2.
\]
\end{proposition}

\begin{proof}
Standard; see Has'mi\u{n}ski\u{\i}~\cite{Hasminski}, Chapter~4.
\end{proof}

\begin{definition}[Embodiment threshold]\label{def:threshold}
The \emph{embodiment threshold} is
$\tau\coloneqq\sup\{\sigma_0:\cL V(x)<0\text{ whenever }
\norm{\sigma(x)}<\sigma_0\}=\sqrt{c/\kappa}$.
Below $\tau$: noise reinforces the orbit. Above $\tau$: noise
disrupts it.
\end{definition}

\begin{definition}[Canonical invariant triple]\label{def:triple}
$\iota(\dm)\coloneqq(T^*,\mu_{\max},\tau)\in
\R_{>0}\times\R_{<0}\times\R_{>0}$.
\end{definition}

\subsection{Proof of Theorem~A}

\begin{proof}[Proof of Theorem~\ref{thm:A}]
Claims (i)--(ii): Axioms~\ref{ax:orbit}, \ref{ax:closure},
\ref{ax:noncollapse}, \ref{ax:hyperbolic} give a compact connected
hyperbolic limit cycle with $T^*>0$. Proposition~\ref{prop:orbital}
gives exponential stability with $c=\abs{\mu_{\max}}$.
Claim~(iii): Proposition~\ref{prop:stoch} with $\tau=\sqrt{c/\kappa}$.
Claim~(iv): Proposition~\ref{prop:winding} in
Section~\ref{sec:contact}.
\end{proof}

%%=============================================================
\section{The operator algebra}\label{sec:operators}
%%=============================================================

Throughout, $\dm=(S,g,Y,\Orb,V,\sigma)$ is a fixed dm3 system
satisfying Axioms~\ref{ax:orbit}--\ref{ax:hyperbolic}.

\subsection{\texorpdfstring{$g$}{g}-series operators: expansion}
\label{subsec:g}

\subsubsection{\texorpdfstring{$g_1$}{g1}: local basin expansion}

\begin{definition}[$g_1$]\label{def:g1}
Let $\rho\colon S\to[0,1]$ be a smooth cutoff with $\rho=0$ on
$\cN_\delta(\Orb)$ and $\rho=1$ outside $\cN_{2\delta}(\Orb)$,
and let $h\colon\R_{\ge 0}\to\R_{\ge 0}$ be smooth with
$\supp h\subset(\delta,R)$ for $R>2\delta$. Define
\[
  Z(x)\coloneqq -\rho(x)\,h(V(x))\,\nabla V(x),\quad
  g_1\coloneqq\exp(Z)=\psi^1,
\]
where $\psi^t$ is the flow of $Z$.
\end{definition}

\begin{proposition}[$\Orb$-invariance of $g_1$]\label{prop:g1inv}
$g_1(\Orb)=\Orb$.
\end{proposition}
\begin{proof}
For $x\in\Orb$, $V(x)=0$, so $\rho(x)=0$ and $Z(x)=0$.
\end{proof}

\begin{proposition}[Basin expansion]\label{prop:g1expand}
The pushforward $Y^{(1)}\coloneqq(g_1)_*Y$ has Lyapunov function
$V^{(1)}=V\circ g_1^{-1}$ satisfying $\dot V^{(1)}\le -c^{(1)}V^{(1)}$
on $\cN^{(1)}(\Orb)\supset\cN(\Orb)$ with $c^{(1)}\ge c$.
The expanded system $\dm^{(1)}=(S,g,Y^{(1)},\Orb,V^{(1)},\sigma)$
is a dm3 system with $\tau^{(1)}\ge\tau$.
\end{proposition}

\begin{conjecture}[Global basin expansion]\label{conj:global}
For appropriate $h$ and $\delta$, the basin of attraction of $\Orb$
under $Y^{(1)}$ is strictly larger than under $Y$.
\end{conjecture}

\subsubsection{\texorpdfstring{$g_2$}{g2}: recursive multi-scale expansion}

\begin{definition}[Phase function and second-scale space]\label{def:phase}
The \emph{phase function} $\varphi\colon\cN(\Orb)\to\R/\mathbb{Z}$
assigns arc-length normalized phase, with $\varphi(p_I)=0$,
$\varphi(p_{VI})=1/4$, $\varphi(p_{III})=1/2$,
$\varphi(p_{VII})=3/4$. Set $S^{(2)}=\R/\mathbb{Z}$ with
projection $\Psi(x)=\varphi(x)$.
\end{definition}

\begin{definition}[$g_2$ and semi-conjugacy]\label{def:g2}
Let $\omega_2>0$, $f\colon S^{(2)}\to\R$ smooth. Define
$Y^{(2)}(\Phi)=\omega_2+f(\Phi)$ with flow $\psi^{(2),t}$, and
$g_2=\psi^{(2),1/\omega_2}$. The orbit $\Orb$ is
\emph{recursively invariant} if $\exists\,\lambda>0$ such that
\[
  \Psi(\phi^t(x)) = \psi^{(2),t/\lambda}(\Psi(x))
  \quad\forall\,x\in\Orb.
\]
Setting $\omega_2=1/T^*$ and $f\equiv 0$ achieves this with $\lambda=1$.
\end{definition}

\begin{conjecture}[Invariant torus stability]\label{conj:torus}
If the semi-conjugacy holds and $Y^{(2)}$ has a hyperbolic limit cycle,
the product system on $S\times S^{(2)}$ admits a stable invariant torus
(KAM-stable in the non-resonant case; mode-locked in the resonant case).
Proved in the toy model in~\cite{Paper2}.
\end{conjecture}

\subsubsection{\texorpdfstring{$g_3$}{g3}: cross-domain expansion and the category \texorpdfstring{$\dmcat$}{dm3}}

\begin{definition}[Category $\dmcat$]\label{def:cat}
The \emph{category} $\dmcat$ has objects: dm3 systems satisfying
Axioms~\ref{ax:orbit}--\ref{ax:hyperbolic}; morphisms: smooth maps
$f\colon S_1\to S_2$ satisfying:
\begin{enumerate}[label=\textup{(M\arabic*)},leftmargin=2em]
  \item $f(\Orb_1)=\Orb_2$,
  \item $f\circ\phi_1^t=\phi_2^{\lambda t}\circ f$ for some $\lambda>0$,
  \item $\alpha\,V_1(x)\le V_2(f(x))\le\beta\,V_1(x)$ for
        $0<\alpha\le\beta<\infty$,
  \item $\norm{\sigma_2(f(x))}\le C\norm{\sigma_1(x)}$ and
        $\tau_2\ge\alpha\tau_1/C$;
\end{enumerate}
composition: function composition; identity: $\mathrm{id}_S$.
\end{definition}

\begin{proposition}[Morphisms preserve dm3 structure]\label{prop:morphism}
If $f\colon\dm_1\to\dm_2$ is a dm3 morphism, then $\dm_2\in\dmcat$.
\end{proposition}
\begin{proof}
Axiom~\ref{ax:orbit}: (M1) and (M2) map phase points to phase points.
Axiom~\ref{ax:closure}: $f(\Orb_1)=\Orb_2$ is invariant by (M2).
Axioms~\ref{ax:primacy}--\ref{ax:response}: (M3) gives a Lyapunov
function for $\Orb_2$ with rate $c'\ge\alpha c/\beta$.
Axiom~\ref{ax:noise}: (M4). Axioms~\ref{ax:noncollapse}--\ref{ax:hyperbolic}:
smoothness of $f$ and (M1)--(M2) with $C^1$-diffeomorphism property
near $\Orb_1$.
\end{proof}

\begin{corollary}[Falsifiability of cross-domain claims]\label{cor:false}
The claim that the dm3 framework applies to a domain $D$ with state
space $S_D$ is equivalent to the existence of a morphism
$f\colon\dm_{\mathrm{base}}\to\dm_D$. This is a falsifiable condition.
\end{corollary}

\subsection{\texorpdfstring{$L$}{L}-operators: coherence and anti-collapse}
\label{subsec:L}

\subsubsection{\texorpdfstring{$L_1$}{L1}: local coherence}

\begin{definition}[$L_1$]\label{def:L1}
With two orbits $\Orb_1,\Orb_2$, inter-orbit distance
$D_{12}(x)=\dg(x,\Orb_1)+\dg(x,\Orb_2)$, and smooth cutoff
$\chi$ supported where $\min\{\dg(x,\Orb_i)\}<\eps$, define
$Z(x)=-\chi(x)\nabla D_{12}(x)$ and $L_1=\exp(Z)$.
\end{definition}

\begin{proposition}[$L_1$ orbit invariance]\label{prop:L1}
$L_1(\Orb_i)=\Orb_i$ for $i=1,2$. There exists $\eps'>\eps$ such
that $\dg(L_1(x),\Orb_i)\ge\dg(x,\Orb_i)$ for
$\dg(x,\Orb_i)\in[\eps,\eps']$.
\end{proposition}
\begin{proof}
$\chi=0$ on each $\Orb_i$, so $Z=0$ there. The gradient push
increases separation in the buffer zone by compactness.
\end{proof}

\begin{conjecture}[$L_1$ hyperbolicity]\label{conj:L1}
For sufficiently small $\norm{Z}_{C^1}$, both $\Orb_i$ remain
hyperbolic under $(L_1)_*Y_i$.
\end{conjecture}

\subsubsection{\texorpdfstring{$L_2$}{L2}: global coherence}

\begin{definition}[$L_2$]\label{def:L2}
The \emph{coherence defect}
$\mathcal{E}(x)=d_{S^{(2)}}(\Psi(\phi^{\Delta t}(x)),
\psi^{(2),\Delta t/\lambda}(\Psi(x)))^2$ measures failure of
semi-conjugacy. Set $\cV(x)=V(x)+\alpha_0\,\mathcal{E}(x)$,
$\alpha_0>0$, and $L_2=\exp(-\nabla\cV)$.
\end{definition}

\begin{proposition}[$L_2$ fixed points]\label{prop:L2}
$L_2(x)=x\iff x\in\Orb$. Moreover $\mathcal{E}|_{\Orb}=0$ and
$\nabla\cV|_{\Orb}=0$.
\end{proposition}
\begin{proof}
$\nabla\cV=0$ iff $\nabla V=0$ (i.e.\ $x\in\Orb$) and
$\nabla\mathcal{E}=0$. On $\Orb$ the semi-conjugacy holds
(Proposition~\ref{def:g2}) so $\mathcal{E}|_{\Orb}=0$ and
$\nabla\mathcal{E}|_{\Orb}=0$.
\end{proof}

\subsubsection{\texorpdfstring{$L_3$}{L3}: anti-collapse}

\begin{definition}[$L_3$]\label{def:L3}
For $\eps_0>0$, define the regularized potential
$U_{12}^\eps(x)=((\dg(x,\Orb_1)^2+\eps_0)^{-1}+(\dg(x,\Orb_2)^2+\eps_0)^{-1})$,
cutoff $\chi_{12}$ in the inter-orbit region,
$Q(x)=\chi_{12}(x)\nabla U_{12}^\eps(x)$,
and $L_3=\exp(\eta Q)$ for small $\eta>0$.
\end{definition}

\begin{proposition}[$L_3$ non-collapse and hyperbolicity]\label{prop:L3}
For sufficiently small $\eta$:
$\dg(L_3(\Orb_1),L_3(\Orb_2))\ge\dg(\Orb_1,\Orb_2)$,
and each $\Orb_i$ remains a hyperbolic limit cycle of $(L_3)_*Y_i$.
\end{proposition}
\begin{proof}
$Q$ points in the direction increasing $\dg(\Orb_1,\Orb_2)$ to first
order. Hyperbolicity is an open $C^1$ condition; $L_3$ is $C^1$-close
to the identity for small $\eta$.
\end{proof}

\subsection{\texorpdfstring{$R$}{R}-operators: resonance}
\label{subsec:R}

\subsubsection{\texorpdfstring{$R_1$}{R1}: detection}

\begin{definition}[$R_1$]\label{def:R1}
Systems $\dm_1,\dm_2$ are in \emph{$k$:$m$ period resonance}
if $kT_1^*=mT_2^*$; in \emph{Lyapunov resonance} if
$\mu_{\max,1}=\mu_{\max,2}$; in \emph{threshold resonance} if
$\tau_1=\tau_2$. Define
\[
  R_1(\dm_1,\dm_2)\coloneqq
  \{(k,m):kT_1^*=mT_2^*\}\cup\mathbf{1}[\mu_{\max,1}=\mu_{\max,2}]
  \cup\mathbf{1}[\tau_1=\tau_2].
\]
$R_1=\emptyset$ iff no resonance; only trivial coproduct
unification is then possible.
\end{definition}

\begin{proposition}[Genericity]\label{prop:R1}
Each resonance locus is codimension-$1$ in $\cP^2$ and has measure
zero. Within each Arnold tongue, $k$:$m$ resonance is structurally
stable under coupling.
\end{proposition}

\subsubsection{\texorpdfstring{$R_2$}{R2}: amplification}

\begin{definition}[$R_2$]\label{def:R2}
Phase mismatch:
$W_{k:m}(\theta_1,\theta_2)=(k\varphi_1(\theta_1)-m\varphi_2(\theta_2))^2$.
Define $Z_{k:m}=-\nabla_{(x_1,x_2)}W_{k:m}$ and
$R_2=\exp(Z_{k:m})$.
\end{definition}

\begin{proposition}[Phase locking and descent]\label{prop:R2}
$\Fix(R_2)\cap(\Orb_1\times\Orb_2)=\{k\varphi_1=m\varphi_2\}$,
consisting of $\lcm(k,m)$ points generically.
Along the flow: $\tfrac{d}{dt}W_{k:m}=-\norm{\nabla W_{k:m}}^2\le 0$.
\end{proposition}
\begin{proof}
Fixed points are zeros of $Z_{k:m}$, i.e.\ critical points of
$W_{k:m}$. On the torus $(\R/\mathbb{Z})^2$, the equation
$k\varphi_1-m\varphi_2=0$ has $\lcm(k,m)$ solutions generically.
Descent: $\dot W_{k:m}=-\norm{\nabla W_{k:m}}^2$ by definition.
\end{proof}

\subsubsection{\texorpdfstring{$R_3$}{R3}: stabilization}

\begin{definition}[$R_3$]\label{def:R3}
Resonance manifold:
$\cM_{k:m}=\{(x_1,x_2)\in\cN(\Orb_1)\times\cN(\Orb_2):
k\varphi_1(x_1)=m\varphi_2(x_2)\}$.
Define $R_3=\exp(-\nabla\dg(\cdot,\cM_{k:m})^2)$.
\end{definition}

\begin{proposition}[Exponential stability]\label{prop:R3}
With $U=d(\cdot,\cM_{k:m})^2$: $\dot U=-4U$ along $R_3$-flow,
so $U(t)=U(0)e^{-4t}$.
Both $\Orb_i$ are fixed by $R_3$.
\end{proposition}
\begin{proof}
$\dot U=-\norm{\nabla U}^2=-4d^2=−4U$ since $\norm{\nabla d}=1$
away from the cut locus. Orbit invariance: $\nabla d(\cdot,\cM_{k:m})$
is tangent to each $\Orb_i$ on $\Orb_i$, so the gradient flow moves
$\Orb_i$ along itself.
\end{proof}

\subsection{\texorpdfstring{$U$}{U}-operators: unification}
\label{subsec:U}

\subsubsection{Admissible spans and pushout existence}

\begin{definition}[Resonant orbit and admissible span]\label{def:span}
$\Orb_\lambda\coloneqq\cM_{k:m}\cap(\Orb_1\times\Orb_2)$
(consists of $\lcm(k,m)$ points). A span
$\dm_1\xleftarrow{p_1}\dm_\lambda\xrightarrow{p_2}\dm_2$
is \emph{admissible} if: (1)~$p_1,p_2$ are dm3 morphisms;
(2)~smooth embeddings; (3)~$\dm_\lambda$ has tubular neighborhoods
in each $S_i$; (4)~$V_1\circ p_1=V_2\circ p_2$ on $\dm_\lambda$.
\end{definition}

\begin{proposition}[Projections are morphisms]\label{prop:proj}
The projections $p_i\colon\cM_{k:m}\to S_i$ satisfy (M1)--(M4)
with $p_i(\Orb_\lambda)=\Orb_i$.
\end{proposition}
\begin{proof}
(M1): $p_i(\Orb_\lambda)=\Orb_i$ by definition.
(M2): $k$:$m$ resonance gives semi-conjugacy with $\lambda=k/m$.
(M3): $V_i\circ p_i$ equivalent near $\Orb_\lambda$ by transversality
of $\cM_{k:m}$ to $V_i$-level sets.
(M4): bounded noise coefficients.
\end{proof}

\begin{proposition}[Pushout existence]\label{prop:pushout}
If the span is admissible, the pushout
$\dm_{12}=\dm_1\sqcup_{\dm_\lambda}\dm_2$ exists in $\dmcat$,
unique up to dm3-isomorphism.
\end{proposition}
\begin{proof}
Conditions (1)--(3) give smooth gluing of $S_1$,$S_2$ along
$p_i(\cM_{k:m})$ by the smooth gluing theorem~\cite{Lee};
(4) gives smooth $V_{12}$. Hyperbolicity of $\Orb_{12}$ follows
from hyperbolicity of $\Orb_1,\Orb_2$ and transversality of the
gluing; uniqueness is categorical.
\end{proof}

\subsubsection{\texorpdfstring{$U_1$, $U_2$, $U_3$}{U1, U2, U3}}

\begin{definition}[$U_1$, $U_2$, $U_3$]\label{def:U}
$U_1(\dm_1,\dm_2)\coloneqq\dm_1\sqcup_{\dm_\lambda}\dm_2$.

The \emph{resonance correspondence}
$\cC_{12}=\{(x_1,x_2)\in S_1\times S_2:(x_1,x_2)\in\cM_{k:m}\}$
defines $U_2(x_1)=\pi_2(\pi_1^{-1}(x_1)\cap\cC_{12})$.

The \emph{synthesized orbit}
$\Orb_{\mathrm{syn}}=\omega\text{-limit set of generic trajectories}$;
$U_3=\exp(-\nabla\dg(\cdot,\Orb_{\mathrm{syn}})^2)$.
\end{definition}

\begin{proposition}[dm3 closure under $U_3$]\label{prop:closure}
After applying $U_3$, the system
$\dm_{12}=(S_{12},g_{12},Y_{12},\Orb_{\mathrm{syn}},V_{12},\sigma_{12})$
satisfies Axioms~\ref{ax:orbit}--\ref{ax:hyperbolic}.
\end{proposition}
\begin{proof}
$\Orb_{\mathrm{syn}}$ is a limit cycle by the $k$:$m$ resonance
collapsing the product torus; hyperbolicity from normal hyperbolicity
of $\cM_{k:m}$ (Proposition~\ref{prop:R3}). Axioms~\ref{ax:orbit}--\ref{ax:response}
hold by construction of $V_{12}$; Axiom~\ref{ax:noise} inherited
with $\tau_{12}\le\min(\tau_i)$; Axioms~\ref{ax:noncollapse}--\ref{ax:hyperbolic}
from the above.
\end{proof}

\subsection{Proof of Theorem~B}

\begin{proof}[Proof of Theorem~\ref{thm:B}]
(i) Well-definedness: $R_1\ne\emptyset$ (resonance assumed);
$R_2$ drives to phase locking (Proposition~\ref{prop:R2});
$R_3$ stabilizes $\cM_{k:m}$ (Proposition~\ref{prop:R3});
admissibility gives the pushout (Proposition~\ref{prop:pushout});
$U_1$--$U_3$ proceed as in Definition~\ref{def:U};
$\dm_{12}\in\dmcat$ by Proposition~\ref{prop:closure}.
(ii)--(iii) Period: $T_{12}^*=kT_1^*/\gcd(k,m)$ (smallest $T>0$
with $T/T_i^*\in\mathbb{Z}$). Stability: slowest contraction dominates.
Threshold: $c_{12}=\min(c_i)$, $\kappa_{12}=\max(\kappa_i)$,
giving $\tau_{12}=\sqrt{c_{12}/\kappa_{12}}\le\min(\tau_i)$.
\end{proof}

%%=============================================================
\section{Boundary operators and constraint geometry}\label{sec:boundary}
%%=============================================================

\subsection{\texorpdfstring{$B_1$}{B1}: identity boundary}

\begin{definition}[$B_1$]\label{def:B1}
$B_1\coloneqq\mathcal{W}^s(\Lambda)=\partial\cB(\Orb)$,
the stable manifold of the boundary invariant set $\Lambda$.
The critical Lyapunov level
$c^*\coloneqq\sup\{c:\{V<c\}\subset\cB(\Orb)\}$
gives inner approximation $\{V=c^*\}\subseteq B_1$.
\end{definition}

\begin{proposition}[$B_1$ invariance]\label{prop:B1}
$\phi^t(B_1)=B_1$ for all $t\in\R$.
\end{proposition}
\begin{proof}
$B_1=\mathcal{W}^s(\Lambda)$ is invariant by definition of stable manifolds.
\end{proof}

\subsection{\texorpdfstring{$B_2$}{B2}: transition boundary}

\begin{definition}[$B_2$ and Arnold tongues]\label{def:B2}
Parameter space: $\cP=\R_{>0}\times\R_{<0}\times\R_{>0}$.
Transition boundary in $\cP^2$:
\[
  B_2=\bigcup_{k,m}\{kT_1^*=mT_2^*\}
  \cup\{\mu_{\max,1}=\mu_{\max,2}\}\cup\{\tau_1=\tau_2\}.
\]
Arnold tongue:
$\mathcal{A}_{k:m}=\{\abs{kT_1^*-mT_2^*}<\eps_{k:m}(A)\}$.
Inside $\mathcal{A}_{k:m}$: $R$-operators valid. Outside: invalid.
\end{definition}

\subsection{\texorpdfstring{$B_3$}{B3}: collapse boundary}

\begin{definition}[$B_3$]\label{def:B3}
$U_{12}^\eps(x)=(\dg(x,\Orb_1)^2+\eps_0)^{-1}
+(\dg(x,\Orb_2)^2+\eps_0)^{-1}$. Define
$C_{12}^*=\sup\{c:\dot U_{12}^\eps>0\text{ whenever }U_{12}^\eps>c\}$
and $B_3=\{U_{12}^\eps=C_{12}^*\}$.
\end{definition}

\begin{proposition}[Smoothness of $B_3$]\label{prop:B3}
For generic $\eps_0$, $B_3$ is a smooth codimension-$1$ submanifold.
\end{proposition}
\begin{proof}
$U_{12}^\eps$ is smooth; Sard's theorem gives generic regular values;
the regular level set theorem applies.
\end{proof}

\subsection{Admissible regions and operator grammar}

\begin{definition}[Admissible regions]\label{def:admissible}
$\cD_S\coloneqq S\setminus(B_1\cup B_3)$;\quad
$\cD_\cP\coloneqq\bigcup_{k,m}\mathcal{A}_{k:m}$.
A configuration $(\dm_1,\dm_2)$ in states $(x_1,x_2)$ is
\emph{fully admissible} if
$(x_i)\in\cD_{S_i}$ and $(\iota(\dm_i))\in\cD_\cP$.
\end{definition}

\begin{theorem}[Operator grammar]\label{thm:grammar}
Let $(\dm_1,\dm_2)$ be fully admissible. The sequence
$L_3\to L_1\to g_1\to L_2\to g_2\to R_1\to R_2\to R_3\to U_1\to U_2\to U_3$
is well-defined and produces a dm3 system at every stage.
Ordering constraints:
(1)~$L_3$ before $g_1$ (separation before expansion);
(2)~$L_1$ before $L_2$ (local before global coherence);
(3)~$g_1$ before $g_2$ (base before recursive expansion);
(4)~$R_1$ before $U_1$ (resonance needed for pushout);
(5)~$R_3$ before $U_3$ (resonance stabilized before synthesis).
Steps involving $L_1$ and $g_2$ are conditional on
Conjectures~\ref{conj:L1} and~\ref{conj:torus}.
The reduced sequence $L_3\to g_1\to R_1\to R_2\to R_3\to U_1\to U_2\to U_3$
holds unconditionally.
\end{theorem}

\begin{proof}
dm3 membership at each stage:
After $L_3$: Proposition~\ref{prop:L3}.
After $L_1$: Proposition~\ref{prop:L1} (conditional).
After $g_1$: Proposition~\ref{prop:g1expand}.
After $L_2$: Proposition~\ref{prop:L2}.
After $g_2$: semi-conjugacy preserved (conditional on Conjecture~\ref{conj:torus}).
After $R_1$: detection only.
After $R_2$, $R_3$: Propositions~\ref{prop:R2},~\ref{prop:R3}.
After $U_1$: Proposition~\ref{prop:pushout}.
After $U_2$: correspondence, no dm3 alteration.
After $U_3$: Proposition~\ref{prop:closure}.
Admissibility maintained since $B_1,B_3$ persist under controlled
operator perturbations (shown case by case above).
\end{proof}

%%=============================================================
\section{Contact geometric physics}\label{sec:contact}
%%=============================================================

\subsection{Why contact geometry}

Seven structural features of dm3 systems force the contact framework:
(1)~hyperbolic limit cycle attractors; (2)~strictly decreasing Lyapunov
functions; (3)~gradient-like operators; (4)~noise-as-fuel stochastic
stability; (5)~embodiment threshold as bifurcation parameter;
(6)~Arnold tongue resonance; (7)~categorical pushouts of dissipative
systems. Features (1)--(2) exclude Hamiltonian mechanics (Liouville's
theorem). Features (4)--(5) require a dissipation parameter evolving
along trajectories. Contact geometry is the unique framework satisfying
all seven simultaneously.

\subsection{The contact manifold}

\begin{definition}[Extended state space and contact form]\label{def:contact-mfld}
Define $M=S\times\R$ with coordinates $(s,z)$, where $z$ is the
\emph{action variable}. Let $\lambda\in\Omega^1(S)$ with $d\lambda=\omega$
a closed non-degenerate 2-form. Set $\alpha=dz-\lambda$.
\end{definition}

\begin{proposition}[Contact structure and Reeb field]\label{prop:contact}
$(M,\alpha)$ is a contact manifold with $R_\alpha=\partial_z$.
\end{proposition}
\begin{proof}
$\alpha\wedge(d\alpha)^n=(dz-\lambda)\wedge(-\omega)^n\ne 0$ since
$\omega^n\ne 0$ and $dz$ is independent of $S$-forms. For Reeb:
$\alpha(\partial_z)=1$ and $\iota_{\partial_z}(-\omega)=0$.
\end{proof}

\subsection{Winding integral and Theorem~A claim (iv)}

\begin{proposition}[Topological invariance of $\oint_\Orb\lambda$]
\label{prop:winding}
\begin{enumerate}[label=\textup{(\roman*)}]
  \item $\oint_\Orb\lambda$ is independent of the choice of $\lambda$
        with $d\lambda=\omega$.
  \item It is invariant under admissible operators fixing $\Orb$ or
        mapping it to a homologous cycle.
  \item It is nonzero when $\Orb$ bounds no chain in $S$.
\end{enumerate}
\end{proposition}
\begin{proof}
(i) $\oint_\Orb df=0$ for any $f$ since $\Orb$ is closed.
(ii) A contactomorphism changes $\lambda$ by an exact form along $\Orb$.
(iii) Stokes' theorem.
\end{proof}

\subsection{Contact Hamiltonian and dynamics}

\begin{definition}[Generative Hamiltonian and equations]\label{def:H}
$H(s,z)=H_{\mathrm{phase}}(s)+H_{\mathrm{diss}}(s,z)$,
where $H_{\mathrm{phase}}=\tfrac12\omega(Y,Y)$ and
$H_{\mathrm{diss}}=-\gamma V(s)e^{-\beta z}$ for $\gamma,\beta>0$.
The full dynamics are $\dot x=X_H(x)+X_{\cR}(x)$, where
$\cR=\tfrac{c}{2}\norm{\dot s-Y(s)}_g^2$ is the Rayleigh function
and $X_{\cR}(s)=-c(\dot s-Y(s))$.
\end{definition}

\begin{proposition}[Recovery of dm3 dynamics]\label{prop:recovery}
$\pi_*X_H(s)=Y(s)-\gamma e^{-\beta z}\nabla V(s)+O(V)$.
On $\Orb$: $\pi_*X_H=Y$.
Along trajectories: $\dot V\le -c\norm{\nabla V}^2\le 0$,
with equality iff $s\in\Orb$.
\end{proposition}

\subsection{Action functional and Euler--Lagrange equations}

\begin{definition}[Generative action]\label{def:action}
$\cS[\gamma]=\int_0^T(\dot z-\lambda(\dot s)-H(s,z))\,dt$
for $\gamma(t)=(s(t),z(t))\colon[0,T]\to M$.
\end{definition}

\begin{theorem}[Contact Euler--Lagrange equations]\label{thm:EL}
Critical points of $\cS$ satisfy
\[
  \dot s=\frac{\partial H}{\partial p}+X_{\cR},\quad
  \dot z=\lambda(\dot s)+H,\quad
  \dot p=-\frac{\partial H}{\partial s}+p\frac{\partial H}{\partial z}.
\]
The contact term $p\,\partial H/\partial z$ is absent from standard
Hamiltonian mechanics.
\end{theorem}

\begin{proposition}[On-orbit action]\label{prop:orbit-action}
$\cS[\gamma|_\Orb]=\oint_\Orb\lambda$ (the topological invariant
of Proposition~\ref{prop:winding}).
\end{proposition}

\subsection{Contact Noether theorem and winding number conservation}

\begin{theorem}[Contact Noether theorem]\label{thm:noether}
If $\{\Phi_\eps\}$ is a one-parameter family of contactomorphisms
with generator $W$ satisfying $\mathcal{L}_W\alpha=f\alpha$ and
$\mathcal{L}_W H=fH$, then $J_W=\alpha(W)$ satisfies
$\dot J_W=f\cdot J_W$ along $X_H$-trajectories.
For strict contactomorphisms ($f=0$): $J_W$ is conserved.
\end{theorem}

\begin{corollary}[Winding number conservation]\label{cor:winding}
Phase rotation on $\Orb$ is a strict contactomorphism. Its conserved
quantity is $\oint_\Orb\lambda$---the topological invariant making
$\Orb$ structurally stable.
\end{corollary}

\subsection{Boundaries as contact submanifolds}

\begin{proposition}[$B_1$ and $B_3$ in contact geometry]\label{prop:contact-bdy}
$\hat B_1=B_1\times\R\subset M$ is a contact hypersurface
($\alpha|_{\hat B_1}$ is a contact form).
$\hat B_3\subset M$ is a Legendrian submanifold
($\alpha|_{\hat B_3}=0$, since $H|_{B_3}=0$).
\end{proposition}

\subsection{Proof of Theorem~C}

\begin{proof}[Proof of Theorem~\ref{thm:C}]
\emph{Step 1 (Floquet coordinates).}
By Floquet theory~\cite{GH}, there exist coordinates $(\rho,\theta)$
near $\Orb$ in $S$ with $\dot\rho=\mu_{\max}\rho+O(\rho^2)$ and
$\dot\theta=\omega+O(\rho)$.

\emph{Step 2 (Contact extension).}
Extend to $M=S\times\R$ with $H$ from Definition~\ref{def:H}.
The contact equations via Theorem~\ref{thm:EL} yield:
$\dot\rho=\mu_{\max}(1-e^{-\beta z})\rho+O(\rho^2)$,
$\dot\theta=\omega+O(\rho)$,
$\dot z=\omega-\abs{\mu_{\max}}\rho^2 e^{-\beta z}+O(\rho^3)$.

\emph{Step 3 (Contactomorphism).}
The coordinate change from the original dm3 coordinates is a
contactomorphism: it preserves $\alpha$ up to a nonvanishing
factor $e^f$, preserving all contact structure.

\emph{Step 4 (Uniqueness).}
The normal form depends only on $(\mu_{\max},\omega,\beta)=\iota(\dm)$.
Any two dm3 systems with the same canonical triple are locally
contact-diffeomorphic.
\end{proof}

%%=============================================================
\section{Structural stability}\label{sec:stability}
%%=============================================================

\subsection{Single-operator and boundary persistence}

\begin{proposition}[Persistence under operators]\label{prop:persist}
For sufficiently small operator parameters:
$g_i(\dm)\in\dmcat$;
$L_i(\dm)\in\dmcat$;
$R_i(\dm_1,\dm_2)\in\dmcat\times\dmcat$;
$U_i(\dm_1,\dm_2)\in\dmcat$.
\end{proposition}
\begin{proof}
Each operator is $C^1$-close to the identity or a dm3 morphism
near $\Orb$; hyperbolicity (open $C^1$ condition) and Lyapunov
descent (continuous in $C^1$) are preserved. See
Propositions~\ref{prop:g1expand}, \ref{prop:L3}, \ref{prop:R3},
\ref{prop:closure}.
\end{proof}

\begin{proposition}[Boundary persistence]\label{prop:bdy-persist}
Under sufficiently small $C^1$ perturbations: $B_1'$, $B_2'$, $B_3'$
exist and vary $C^1$-smoothly.
\end{proposition}
\begin{proof}
$B_1$: stable manifold theorem applied to perturbed $\Lambda'$.
$B_2$: implicit function theorem on $kT_1^*-mT_2^*=0$.
$B_3$: regular level set theorem on $U_{12}^\eps$.
\end{proof}

\subsection{Proof of Theorem~D}

\begin{proof}[Proof of Theorem~\ref{thm:D}]
\emph{Step 1.} Each operator produces a $C^1$-perturbation of $Y$
of size $\delta_i$ (bounded by its parameter; see
Definition~\ref{def:g1} for $g_1$: $O(\delta)$; $L_1,L_3$: $O(\eps)$,
$O(\eta)$; $R_2,R_3$: $O(\norm{W_{k:m}}_{C^1})$).

\emph{Step 2.} The composed perturbation satisfies
$\norm{Y_N-Y}_{C^1}\le\prod_{i=1}^N(1+\delta_i)-1$.
Choosing parameters so this product satisfies
$\prod(1+\delta_i)-1<\varepsilon_0$ keeps the composition within
the $\varepsilon_0$-ball.

\emph{Step 3.} Within the $\varepsilon_0$-ball, Shub's theorem~\cite{Shub}
gives a unique hyperbolic $\Orb'$ near $\Orb$ with
$d_{C^1}(\Orb',\Orb)=O(\varepsilon_0)$. The winding number is preserved
by Proposition~\ref{prop:winding}.

\emph{Step 4.} $T^{*\prime}$, $\mu_{\max}'$, $\tau'$ vary by
$O(\varepsilon_0)$ (period: implicit function theorem; Floquet:
spectral perturbation; threshold: continuity of $\sqrt{c/\kappa}$).

\emph{Step 5.} Proposition~\ref{prop:bdy-persist}.

\emph{Step 6.} $\varepsilon_0<\varepsilon_0$ is satisfied throughout,
so Axioms~\ref{ax:orbit}--\ref{ax:hyperbolic} hold for $\dm'$
by Steps~3--5.
\end{proof}

\begin{remark}[Stability radius]\label{rem:radius}
$\varepsilon_0=\dfrac{\abs{\mu_{\max}}}{2(1+\sup_\Orb\norm{\Hess V}_\infty)}$
(equal to $\tfrac{1}{3}$ in the toy model with $\mu_{\max}=-2$,
$\sup_\Orb\norm{\Hess V}_\infty=2$):
larger $\abs{\mu_{\max}}$ (faster contraction) and flatter $V$ both
enlarge the stability radius. This bound is computable from
$\iota(\dm)$ and the Lyapunov geometry.
\end{remark}

%%=============================================================
\section{Discussion}\label{sec:discussion}
%%=============================================================

\subsection{Summary of main results}

This paper develops a geometric framework for dissipative dynamical
systems with structured limit cycles. The dm3 system, formalized
through eight axioms, is the central object. An operator algebra
acting on the category $\dmcat$ governs expansion ($g$-operators),
coherence ($L$-operators), resonance ($R$-operators), and unification
($U$-operators). The entire structure is embedded in contact geometry
on $M=S\times\R$.

The four main theorems establish: (A)~existence and exponential
stability of the dm3 orbit, with topological invariant
$\oint_\Orb\lambda$ and stochastic stability below $\tau$;
(B)~closure of $\dmcat$ under $\cU$, with the prediction
$\tau_{12}\le\min(\tau_i)$; (C)~the universal contact normal form
with parameters $(\mu_{\max},\omega,\beta)$; and (D)~$C^1$-structural
stability with explicit radius $\varepsilon_0$.

\subsection{Relation to nonequilibrium thermodynamics}

The action variable $z$ plays the role of entropy: it increases
monotonically ($\dot z>0$ on $\Orb$) and the contact form
$\alpha=dz-\lambda$ encodes the first and second laws simultaneously.
The embodiment threshold $\tau$ functions as a noise-induced phase
transition. The result $\tau_{12}\le\min(\tau_i)$ says that
higher-order organization requires lower effective temperature,
consistent with observations in biological and physical oscillator
networks and with stochastic thermodynamic
theory~\cite{Hasminski,PRK}.

\subsection{Relation to coupled oscillator theory}

Period resonance, Arnold tongues, and mode locking are
classical~\cite{PRK,Izhikevich}. The contribution here is embedding
these phenomena in a categorical structure: resonance is the
precondition for the pushout in $\dmcat$. The prediction
$\tau_{12}\le\min(\tau_i)$ is new and testable.

\subsection{Open problems}

\begin{enumerate}[label=(\arabic*),leftmargin=*]
\item \textbf{Global basin expansion} (Conjecture~\ref{conj:global}):
  conditions on $S$ for global (not just local) basin expansion.

\item \textbf{Invariant torus stability} (Conjecture~\ref{conj:torus}):
  KAM theory (quasi-periodic) and Poincar\'e--Melnikov
  (resonant) in the abstract setting. Proved in~\cite{Paper2}
  for the toy model.

\item \textbf{Higher-order resonance and Devil's staircase}:
  measure-theoretic and fractal properties of the full Arnold
  tongue diagram in $\cD_\cP$.

\item \textbf{Infinite-dimensional extensions}: Hilbert/Banach
  manifold generalizations; infinite-dimensional Floquet theory;
  PDE and field-theory dm3 systems.

\item \textbf{Categorical completeness of $\dmcat$}: does $\dmcat$
  have all finite limits? colimits? Is it a topos?
\end{enumerate}

\subsection{Companion paper}

The companion paper~\cite{Paper2} instantiates every definition and
operator on the explicit system
$\dot r=r(1-r^2)+2(r-1)e^{-z}$,
$\dot\theta=1$,
$\dot z=r^2-2(r-1)^2e^{-z}$
on $S=\R^2_{>0}$, verifying all theoretical predictions including
the global attractor, four bifurcations, stationary SDE measure,
and Conjecture~\ref{conj:torus} in the $1$:$2$ resonant case.

%%=============================================================
\appendix
%%=============================================================

\section{The category \texorpdfstring{$\dmcat$}{dm3}: formal development}
\label{app:cat}

[Full categorical development: objects, morphisms, composition law,
associativity, identity, pushouts (proved in
Proposition~\ref{prop:pushout}), coproducts (disjoint union when
$R_1=\emptyset$), and initial/terminal objects if they exist.
Approximately 3--4 pages.]

\section{Supporting lemmas}\label{app:lemmas}

[Proofs deferred from the main text. Approximately 2--3 pages.]

%%=============================================================
\section*{Acknowledgments}
The author acknowledges with gratitude the foundational influence
of the teachings of Paramahamsa Nithyananda and the yogic
scriptural tradition from which this work derives. The mathematical
formalization presented here is understood by the author as a
translation of principles encountered through that tradition.

\begin{thebibliography}{99}
%%=============================================================

\bibitem{Arnold}
V.~I.~Arnold,
\emph{Mathematical Methods of Classical Mechanics},
2nd ed., Springer, New York, 1989.

\bibitem{Bravetti2017}
A.~Bravetti,
\emph{Contact Hamiltonian mechanics},
Ann.\ Phys.\ \textbf{376} (2017), 17--39.

\bibitem{deLeonLainz2019}
M.~de~Le{\'o}n and M.~Lainz~Valc{\'a}zar,
\emph{Contact Hamiltonian systems},
J.\ Math.\ Phys.\ \textbf{60} (2019), 102902.

\bibitem{Gaset2020}
J.~Gaset, X.~Gr{\`a}cia, M.~C.~Mu{\~n}oz-Lecanda, X.~Rivas,
and N.~Rom{\'a}n-Roy,
\emph{New contributions to the Hamiltonian and Lagrangian contact
formalisms for dissipative mechanical systems and optimal control},
J.~Geom.\ Mech.\ \textbf{12} (2020), 543--574.

\bibitem{Geiges}
H.~Geiges,
\emph{An Introduction to Contact Topology},
Cambridge University Press, Cambridge, 2008.

\bibitem{GH}
J.~Guckenheimer and P.~Holmes,
\emph{Nonlinear Oscillations, Dynamical Systems, and Bifurcations
of Vector Fields},
Springer, New York, 1983.

\bibitem{Hasminski}
R.~Z.~Has{'}mi{\u{n}}ski{\u{\i}},
\emph{Stochastic Stability of Differential Equations},
Sijthoff \& Noordhoff, Alphen aan den Rijn, 1980.

\bibitem{HPS}
M.~W.~Hirsch, C.~C.~Pugh, and M.~Shub,
\emph{Invariant Manifolds},
Lecture Notes in Mathematics 583, Springer, Berlin, 1977.

\bibitem{Izhikevich}
E.~M.~Izhikevich,
\emph{Dynamical Systems in Neuroscience},
MIT Press, Cambridge, MA, 2007.

\bibitem{Lee}
J.~M.~Lee,
\emph{Introduction to Smooth Manifolds},
2nd ed., Springer, New York, 2013.

\bibitem{Morrison1986}
P.~J.~Morrison,
\emph{A paradigm for joined Hamiltonian and dissipative systems},
Physica D \textbf{18} (1986), 410--419.

\bibitem{PRK}
A.~Pikovsky, M.~Rosenblum, and J.~Kurths,
\emph{Synchronization: A Universal Concept in Nonlinear Sciences},
Cambridge University Press, Cambridge, 2001.

\bibitem{Shub}
M.~Shub,
\emph{Global Stability of Dynamical Systems},
Springer, New York, 1987.

\bibitem{Paper2}
Pablo Nogueira Grossi,
\emph{The dm3 Operator: Explicit Toy Model and Global Dynamical
Analysis},
Appendix~A of this volume. Zenodo: \texttt{10.5281/zenodo.19117400}.

\end{thebibliography}
\cleardoublepage


%% ════════════════════════════════════════════════════════════
\part{Biological Instantiations}
%% ════════════════════════════════════════════════════════════

%% ============================================================

\title[dm3 Operator: Toy Model]{%
  The dm3 Operator: Explicit Toy Model\\
  and Global Dynamical Analysis}

\author{Pablo Nogueira Grossi}
\address{Academic Coordinator, UCEDA School (ESL), Elizabeth, NJ\\
Newark, NJ, USA}
\email{pablogrossi@hotmail.com}

\keywords{contact geometry, limit cycles, bifurcation theory,
  invariant measures, structural stability, resonance, normal form,
  stochastic stability, global attractor}

\subjclass[2020]{37C10, 37C27, 37G15, 37H10, 53D10, 60H10}



}

%%=============================================================
\section{Introduction}\label{sec:intro}
%%=============================================================

\subsection{Purpose and scope}

This paper is the companion to~\cite{Paper1}, which develops the
generative contact mechanics framework in full generality. The purpose
here is verification: we show that every abstract definition, operator,
and theorem of~\cite{Paper1} is instantiated by an explicit, computable
dynamical system, and that all theoretical predictions can be checked
by direct computation or straightforward analysis.

The explicit system is chosen for concreteness, not generality.
It is the simplest system exhibiting all eight axioms of~\cite{Paper1},
the full operator algebra, the contact geometric structure, and the
global dynamics that the abstract theory predicts.

\subsection{The system}

The base manifold is $S=\R^2_{>0}$ with polar coordinates $(r,\theta)$,
$r>0$, $\theta\in\R/2\pi\mathbb{Z}$. The contact manifold is
$M=S\times\R$ with coordinates $(r,\theta,z)$. The defining equations are
\begin{align}
  \dot r &= r(1-r^2) + 2(r-1)e^{-z},\label{eq:rdot}\\
  \dot\theta &= 1 - \frac{2(r-1)}{r}e^{-z},\label{eq:thetadot}\\
  \dot z &= r^2 - 2(r-1)^2 e^{-z},\label{eq:zdot}
\end{align}
with contact form $\alpha=dz-r^2\,d\theta$.

\subsection{Main results}

\begin{theorem}[A: global attractor]\label{thm:A}
The global attractor of the system~\eqref{eq:rdot}--\eqref{eq:zdot}
restricted to $\cB(\Orb)\times\R_{>0}$ is
\[
  \mathcal{A}_{\mathrm{full}} = \Orb_{12},
\]
the $1$:$2$ resonant orbit on $\Orb\times\Orb_2$. Every trajectory
in the fully admissible region converges to $\Orb_{12}$.
\end{theorem}

\begin{theorem}[B: invariant torus conjecture]\label{thm:B}
The coupled system on $S\times S^{(2)}$ with the $g_2$ semi-conjugacy
admits a normally hyperbolic invariant circle $\Orb_{12}$ in the
$1$:$2$ resonant case. The transverse Lyapunov exponent is $-3$,
and $\Orb_{12}$ is exponentially stable. This proves
Conjecture~4.15 of~\cite{Paper1} in the toy model.
\end{theorem}

\begin{theorem}[C: bifurcation diagram]\label{thm:C}
As parameters $(\gamma,\beta,\eta)$ vary from their baseline values
$(\gamma,\beta,c,\eta)=(2,1,2,1)$:
\begin{enumerate}[label=\textup{(\roman*)}]
  \item \textup{(Contact Hopf)} At $\gamma=\gamma^*=e^{z_0}$, the
        system undergoes a supercritical contact Hopf bifurcation.
  \item \textup{(Saddle-node)} At $\eta=\eta^*\approx 0.15$, two
        limit cycles collide and the multi-orbit structure collapses.
  \item \textup{(Neimark--Sacker)} At detuning $\abs{\Delta}=\Delta^*$,
        the resonant orbit $\Orb_{12}$ loses stability and a 2-torus
        bifurcates. This boundary is $\partial\mathcal{A}_{1:2}=B_2$.
  \item \textup{(Slow-fast crossover)} At $\beta=\beta^*$, the
        contact correction crosses over from dominant to negligible;
        this is a smooth parameter-space transition, not a sharp
        bifurcation.
\end{enumerate}
\end{theorem}

\begin{theorem}[D: stationary measure]\label{thm:D}
For the SDE extension with noise $\sigma>0$, the stationary measure
$\mu_\sigma$ satisfies:
\[
  \mu_\sigma(\cN_\delta(\Orb)) = \erf\!\left(\frac{\delta}{\sigma\sqrt{2}}\right).
\]
For $\sigma<\tau$: $\mu_\sigma$ concentrates on $\Orb$ exponentially
fast as $\sigma\to 0$. For $\sigma>\tau$: $\mu_\sigma$ spreads to
fill $\cB(\Orb)$.
\end{theorem}

\subsection{Organization}

Section~\ref{sec:setup} defines the system and verifies the eight
axioms.
Section~\ref{sec:operators} instantiates all operators explicitly.
Section~\ref{sec:invariants} identifies the invariant sets and
classifies their stability.
Section~\ref{sec:portrait} gives the global phase portrait and
proves Theorem~A.
Section~\ref{sec:invariants} proves Theorem~B (invariant torus conjecture).
Section~\ref{sec:bifurcations} proves Theorem~C (bifurcation diagram).
Section~\ref{sec:measures} proves Theorem~D (stationary measures).
Section~\ref{sec:normal} verifies the contact normal form of
Theorem~C of~\cite{Paper1}.
Section~\ref{sec:discussion} discusses the results and their
implications for~\cite{Paper1}.

%%=============================================================
\section{The toy model: setup and axiom verification}\label{sec:setup}
%%=============================================================

\subsection{The dm3 system}

\begin{definition}[Toy model dm3 system]\label{def:toy}
The toy model dm3 system is
$\dm_0=(S,g,Y,\Orb,V,\sigma)$ where:
\begin{itemize}[leftmargin=*]
  \item $S=\R^2_{>0}$ with polar coordinates $(r,\theta)$ and
        Euclidean metric $g$;
  \item $Y(r,\theta)=r(1-r^2)\partial_r+\partial_\theta$;
  \item $\Orb=\{r=1\}$ with minimal period $T^*=2\pi$;
  \item $V(r,\theta)=(r-1)^2$;
  \item $\sigma(r,\theta)=\sigma_0\partial_r$ for constant $\sigma_0>0$.
\end{itemize}
\end{definition}

\subsection{Explicit flow}

\begin{proposition}[Explicit flow formula]\label{prop:flow}
The flow of $Y$ is
\[
  \phi^t(r_0,\theta_0) =
  \left(\frac{1}{\sqrt{1+(r_0^{-2}-1)e^{-4t}}},\;
  \theta_0 + t\bmod 2\pi\right).
\]
\end{proposition}
\begin{proof}
The angular equation $\dot\theta=1$ integrates to $\theta(t)=\theta_0+t$.
The radial equation $\dot r=r(1-r^2)$ is Bernoulli; the substitution
$u=r^{-2}$ gives $\dot u=-4u+4$, with solution
$u(t)=1+(u_0-1)e^{-4t}$. Inverting gives the stated formula.
\end{proof}

\subsection{Canonical invariant triple}

\begin{proposition}[Canonical invariants]\label{prop:canonical}
The canonical invariant triple of $\dm_0$ is
\[
  \iota(\dm_0) = (T^*,\mu_{\max},\tau) = (2\pi,-2,\tau_0),
\]
where $\tau_0$ is computed in Section~\ref{sec:measures}.
\end{proposition}
\begin{proof}
Period: $T^*=2\pi$ is the period of $\dot\theta=1$ on $\Orb$.
Floquet exponent: linearizing $\dot r=r(1-r^2)$ at $r=1$ with
$r=1+\rho$: $\dot\rho\approx-2\rho$, so $\mu_{\max}=-2$.
\end{proof}

\subsection{Axiom verification}\label{subsec:axioms}

The eight dm3 axioms of~\cite{Paper1} are verified explicitly
for the toy model on $M=\R^2_{>0}\times\R$ with
equations~\eqref{eq:rdot}--\eqref{eq:zdot}. Each axiom is checked
with explicit derivations and cross-references to the global results
of Sections~\ref{sec:invariants}--\ref{sec:measures}.

\begin{proposition}[All eight axioms hold]\label{prop:axioms}
The system $\dm_0$ satisfies Axioms~1--8 of~\cite{Paper1}.
\end{proposition}

\begin{proof}
\textbf{Axiom~1 (Orbit identity).}
The four phase points $p_I=(1,0)$, $p_{VI}=(1,\pi/2)$,
$p_{III}=(1,\pi)$, $p_{VII}=(1,3\pi/2)$ lie on $\Orb=\{r=1\}$
and are cycled by $\phi^{\pi/2}$ (Proposition~\ref{prop:flow}).

\textbf{Axiom~2 (Generative closure).}
On $\Orb$: $\dot r = 1\cdot(1-1)+0=0$, so $r$ remains $1$
under the flow. $\Orb$ is compact, connected, and invariant.

\textbf{Axiom~3 (Structural primacy).}
$V=(r-1)^2$. Computing:
\[
  \dot V = 2(r-1)\dot r
  = 2(r-1)\bigl[r(1-r^2)+2(r-1)e^{-z}\bigr]
  = -2(r-1)^2\bigl[r(1+r)-2e^{-z}\bigr].
\]
For $r$ near $1$ and $z>0$: $r(1+r)\approx 2>2e^{-z}$,
so $\dot V<0$ for $r\ne 1$. Stability is encoded in $Y$ and $V$,
not in external constraints (consistent with Axiom~7).

\textbf{Axiom~4 (Perturbation response).}
Near $\Orb$ with $\rho=r-1$ small, expanding to first order in
$\rho$:
\[
  \dot\rho = -2(1-e^{-z})\rho + O(\rho^2).
\]
For $z>0$: the transverse eigenvalue
$\lambda(z)=-2(1-e^{-z})<0$, so $\dot V\le -c(z)V$ with
$c(z)=4(1-e^{-z})>0$.

\textbf{Axiom~5 (Noise-as-fuel).}
For the It\^{o} SDE $dr=\dot r\,dt+\sigma\,dW_t$, the generator
applied to $V=(r-1)^2$ gives
\[
  \cL V = \dot V + \tfrac{\sigma^2}{2}\cdot 2
  \le -4V(1-e^{-z}) + \sigma^2.
\]
With $c=4(1-e^{-z})$ and $\kappa=1$: asymptotically ($z\to\infty$),
$c\to 4$ and $\tau=\sqrt{c/\kappa}\to 2$.

\textbf{Axiom~6 (Non-collapse).}
$\dot\theta=1\ne 0$ on $\Orb$; minimal period $T^*=2\pi>0$.

\textbf{Axiom~7 (Non-coercion).}
The vector field is $C^\infty$ on $M$; no discontinuous
projection or hard boundary enforcement appears.

\textbf{Axiom~8 (Hyperbolicity).}
From the linearization above: $\mu_{\max}=-2<0$ (asymptotic
rate). All non-trivial Floquet multipliers satisfy $|\mu_j|<1$.
\end{proof}

\begin{lemma}[Transverse stability and embodiment]\label{lem:transverse}
The transverse eigenvalue $\lambda(z)=-2(1-e^{-z})$ satisfies:
\begin{enumerate}[label=\textup{(\roman*)}]
  \item $\lambda(0)=0$: the orbit is neutrally stable at $z=0$,
  \item $\lambda(z)<0$ for all $z>0$: the orbit is exponentially
        stable once action has accumulated,
  \item $\lambda(z)\to -2$ as $z\to\infty$: full dm3 attraction
        is recovered asymptotically.
\end{enumerate}
\end{lemma}

\begin{remark}[Embodiment threshold in the linearization]
The neutral stability at $z=0$ is not a deficiency --- it is
the precise dynamical content of the embodiment threshold.
The orbit earns its stability by accumulating action $z>0$.
This is the contact-geometric mechanism: the action variable $z$
mediates the transition from neutral to exponentially stable
behavior. The asymptotic embodiment threshold is
$\tau=2$ (from $c=4$, $\kappa=1$), and the structural stability
radius is
\[
  \varepsilon_0
  = \frac{|\mu_{\max}|}{2(1+\sup_\Orb\|\Hess V\|_\infty)}
  = \frac{2}{2(1+2)}
  = \frac{1}{3}.
\]
\end{remark}

The canonical invariant triple is therefore
$\iota(\dm_0)=(T^*,\mu_{\max},\tau)=(2\pi,-2,2)$,
consistent with Proposition~\ref{prop:canonical}.

\subsection{The contact manifold}

\begin{proposition}[Contact structure]\label{prop:contactstruct}
On $M=S\times\R$, the form $\alpha=dz-r^2\,d\theta$ satisfies
$\alpha\wedge d\alpha=2r\,dz\wedge dr\wedge d\theta\ne 0$
for $r>0$. The Reeb vector field is $R_\alpha=\partial_z$.
\end{proposition}
\begin{proof}
$d\alpha=-2r\,dr\wedge d\theta$. Then
$\alpha\wedge d\alpha=(dz-r^2\,d\theta)\wedge(-2r\,dr\wedge d\theta)
=-2r\,dz\wedge dr\wedge d\theta\ne 0$ for $r>0$.
$\alpha(\partial_z)=1$ and $\iota_{\partial_z}d\alpha=0$.
\end{proof}

\subsection{Full contact equations of motion}

Setting $\gamma=c=2$ and $\beta=1$, the contact equations are
\begin{align}
  \dot r &= \underbrace{r(1-r^2)}_{\text{dm3}}
           + \underbrace{2(r-1)e^{-z}}_{\text{contact}},
           \label{eq:r-full}\\
  \dot\theta &= 1 - \frac{2(r-1)}{r}e^{-z},\label{eq:theta-full}\\
  \dot z &= r^2 - 2(r-1)^2 e^{-z}.\label{eq:z-full}
\end{align}

\begin{proposition}[Behavior on $\Orb$]\label{prop:onorbit}
On $\Orb$ ($r=1$): $\dot r=0$, $\dot\theta=1$, $\dot z=1$.
The system rotates at unit angular speed and accumulates action
at rate $1$.
\end{proposition}

\begin{proposition}[Embodiment threshold in the toy model]\label{prop:threshold}
Near $\Orb$ with $r=1+\rho$, $\rho$ small:
\[
  \dot\rho \approx -2\rho(1-e^{-z}).
\]
For $z>0$: $\dot\rho<0$ (orbit attracting). For $z=0$: $\dot\rho=0$
(neutral). The embodiment threshold is the first crossing $z>0$.
Larger $z$ gives stronger attraction; this is the mathematical
content of embodiment: the orbit becomes more stable as action
accumulates.
\end{proposition}

%%=============================================================
\section{Operator instantiation}\label{sec:operators}
%%=============================================================

\subsection{\texorpdfstring{$g$}{g}-operators}

\subsubsection{$g_1$: local basin expansion}

Take cutoff $\rho$ supported in $r\in(0.8,1.2)$ and $h(v)=v$.
The modified radial equation is
\[
  \dot r^{(1)} = r(1-r^2) + 2\rho(r)(r-1).
\]
At $r=1$: $\dot r^{(1)}=0$ (orbit invariant). For $r\in(0.8,1.2)$,
the convergence rate to $r=1$ increases:
$\dot V^{(1)}\approx-(4+\eps)V^{(1)}$ for some $\eps>0$.
Enlarged neighborhood: $\cN^{(1)}(\Orb)\supset\cN(\Orb)$. $\checkmark$

\subsubsection{$g_2$: recursive expansion}

Second-scale space $S^{(2)}=\R/\mathbb{Z}$ with coordinate
$\Phi=\theta/2\pi$ (phase normalized to $[0,1)$).
The projection $\Psi(r,\theta)=\theta/2\pi$.

The macro-vector field $Y^{(2)}\equiv 1/T^*=1/2\pi$ gives
$\dot\Phi=1/2\pi$ and macro-period $T^{(2)}=1$.

Semi-conjugacy on $\Orb$: $\Psi(\phi^t(1,\theta_0))=(\theta_0+t)/2\pi
=\Psi(1,\theta_0)+t/2\pi=\psi^{(2),t}(\Psi(1,\theta_0))$.
The semi-conjugacy holds exactly on $\Orb$ with $\lambda=1$. $\checkmark$

\subsubsection{$g_3$: cross-domain morphism}

For a scaled system with $\Orb'=\{r=R\}$, $R>0$:
\[
  g_3(r,\theta) = (Rr,\theta).
\]
Verification of (M1)--(M4):
(M1) $g_3(\{r=1\})=\{r=R\}=\Orb'$.
(M2) $g_3\circ\phi^t=\phi'^t\circ g_3$ with $\lambda=1$ (both
     have unit angular speed).
(M3) $V'\circ g_3=(Rr-R)^2=R^2(r-1)^2=R^2 V$, so $\alpha=\beta=R^2$.
(M4) Radially uniform $\sigma$ gives $C=1$, $\tau'=R^2\tau$. $\checkmark$

\subsection{\texorpdfstring{$L$}{L}-operators}

\subsubsection{$L_1$: local coherence}

For a second orbit $\Orb_2=\{r=R\}$, $R=2$:
\[
  D_{12}(r) = \abs{r-1} + \abs{r-2},\quad
  Z(r) = -\chi(r)\,\mathrm{sgn}(2r-3).
\]
$Z=0$ at $r=1$ and $r=2$ (orbit invariant). In $(1,2)$: $Z$
pushes toward midpoint, increasing the buffer. $\checkmark$

\subsubsection{$L_2$: global coherence}

Coherence defect: $\mathcal{E}(r,\theta)=(\Phi-\theta/2\pi)^2$.
On $\Orb$: $\mathcal{E}=0$ by the $g_2$ semi-conjugacy.

Two-scale Lyapunov: $\mathcal{V}=(r-1)^2+\alpha_0(\Phi-\theta/2\pi)^2$.

$L_2=\exp(-\nabla\mathcal{V})$ drives both to $\Orb$ and to phase
coherence. Fixed points: $r=1$ and $\Phi=\theta/2\pi$, i.e., $\Orb$.
Consistent with Proposition~4.29 of~\cite{Paper1}. $\checkmark$

\subsubsection{$L_3$: anti-collapse}

With $\Orb_1=\{r=1\}$, $\Orb_2=\{r=2\}$, $\eps_0=0.01$:
\[
  U_{12}^\eps(r)
  = \frac{1}{(r-1)^2+0.01} + \frac{1}{(r-2)^2+0.01}.
\]
$\nabla U_{12}^\eps$ computed explicitly; minimum at $r=1.5$
(by symmetry $r\mapsto 3-r$). At $r=1.5$: $U_{12}^\eps\approx 7.69$.
Collapse boundary: $B_3=\{r\approx 1.35\}\cup\{r\approx 1.65\}$. $\checkmark$

\subsection{\texorpdfstring{$R$}{R}-operators}

Add a second dm3 system with $\Orb_2=\{r=1\}$ and period $T_2^*=\pi$
(angular speed $\dot\theta_2=2$). Canonical invariants:
$\iota(\dm_2)=(\pi,-2,\tau_2)$.

\subsubsection{$R_1$: resonance detection}

$F_{1:2}(T_1^*,T_2^*)=1\cdot 2\pi-2\cdot\pi=0$.
The systems are in $1$:$2$ period resonance.
Lyapunov resonance: $\mu_{\max,1}=\mu_{\max,2}=-2$. $\checkmark$

\subsubsection{$R_2$: amplification}

Phase mismatch: $W_{1:2}(\theta_1,\theta_2)=(\theta_1-2\theta_2)^2
\bmod 2\pi$.

Gradient field:
\[
  Z_{1:2} = \begin{pmatrix}-2(\theta_1-2\theta_2)\\
  4(\theta_1-2\theta_2)\end{pmatrix}.
\]
Fixed points: $\theta_1=2\theta_2\bmod 2\pi$ --- two points on
$\Orb_1\times\Orb_2$ (as $\lcm(1,2)=2$). $\checkmark$

\subsubsection{$R_3$: stabilization}

Resonance manifold:
$\cM_{1:2}=\{(\theta_1,\theta_2):\theta_1=2\theta_2\bmod 2\pi\}
\subset\T^2$.

Distance function: $U=(\theta_1-2\theta_2)^2/5$ (normalized).
Generator: $\dot U=-4U$, exponential stability at rate $4$.
Both $\Orb_i$ preserved. $\checkmark$

\subsection{\texorpdfstring{$U$}{U}-operators}

\subsubsection{$U_1$: pushout}

Unified state space: $S_{12}$ obtained by gluing $S_1,S_2$ along
$\cM_{1:2}$.

Unified vector field:
$Y_{12}(\theta_1,\theta_2)=\partial_{\theta_1}+2\partial_{\theta_2}$.

Unified limit cycle:
$\Orb_{12}=\{(\theta_1,\theta_2):\theta_1=2\theta_2\}$, $T_{12}^*=2\pi$.

Canonical invariants:
$T_{12}^*=1\cdot 2\pi/\gcd(1,2)=2\pi$, $\mu_{\max,12}=-2$,
$\tau_{12}=\min(\tau_1,\tau_2)$. $\checkmark$

\subsubsection{$U_2$: translation}

Resonance correspondence: $\cC_{12}=\{(\theta_1,\theta_2):\theta_1=2\theta_2\}$.
$U_2(\theta_1)=\theta_1/2\bmod\pi$: smooth, well-defined map. $\checkmark$

\subsubsection{$U_3$: synthesis}

$\Orb_{\mathrm{syn}}=\Orb_{12}$, the diagonal of $\T^2$.
$U_3=\exp(-\nabla(\theta_1-2\theta_2)^2/5)$ attracts to
$\Orb_{12}$ at rate $4$. $\checkmark$

\subsection{Boundary operators}

\subsubsection{$B_1$: identity boundary}

Basin $\cB(\Orb)=\R^2_{>0}\setminus\{0\}$ (entire punctured plane;
no other attractors). Boundary invariant set: $\Lambda=\{r=0\}$.
$B_1=\{r=0\}$, the repelling singularity. Critical Lyapunov level:
$c^*=+\infty$ ($V$ decreases everywhere). $\checkmark$

\subsubsection{$B_2$: transition boundary}

In the $(T_1^*,T_2^*)$-plane with $T_1^*=2\pi$ fixed, the $1$:$2$
resonance condition $T_2^*=\pi$ is a point. The Arnold tongue
$\mathcal{A}_{1:2}$ is an interval $T_2^*\in(\pi-\eps,\pi+\eps)$
for coupling amplitude-dependent width $\eps=\eps(A)$.
$B_2=\partial\mathcal{A}_{1:2}=\{T_2^*=\pi\pm\eps\}$. $\checkmark$

\subsubsection{$B_3$: collapse boundary}

With $\eps_0=0.01$: $B_3=\{r\approx 1.35\}\cup\{r\approx 1.65\}$
(symmetric about $r=1.5$, computed numerically from
$U_{12}^\eps(r)=C_{12}^*\approx 7.69$). $\checkmark$

%%=============================================================
\section{Invariant sets and stability classification}\label{sec:invariants}
%%=============================================================

\subsection{Identification of invariant sets}

The full system~\eqref{eq:r-full}--\eqref{eq:z-full}
has the following invariant sets:

\begin{enumerate}[label=(\arabic*)]
  \item \textbf{dm3 orbit} $\Orb=\{r=1\}$, a hyperbolic limit cycle.
  \item \textbf{Origin} $\{r=0\}$, a repelling singularity ($B_1$).
  \item \textbf{Resonant orbit} $\Orb_{12}=\{(\theta_1,\theta_2):
        \theta_1=2\theta_2\}\subset\T^2$, the global attractor
        of the product system.
  \item \textbf{Collapse boundary} $B_3$, a normally hyperbolic
        repelling barrier.
\end{enumerate}

\subsection{Stability classification}

\begin{proposition}[dm3 orbit stability]\label{prop:dm3stable}
Linearizing~\eqref{eq:r-full} at $r=1$ with $r=1+\rho$:
\[
  \dot\rho = -2\rho + 2\rho e^{-z} = -2\rho(1-e^{-z}).
\]
For $z>0$: transverse eigenvalue $\lambda_{\dm3}=-2(1-e^{-z})<0$.
The orbit is exponentially stable for all $z>0$.
\end{proposition}

\begin{proposition}[Resonant orbit stability]\label{prop:resstable}
Let $\delta=\theta_1-2\theta_2$ (transverse deviation from $\Orb_{12}$).
Then $\dot\delta=\dot\theta_1-2\dot\theta_2=1-2\cdot 2=-3$.
So $\delta(t)=\delta(0)e^{-3t}$: transverse exponent $-3$,
exponential stability. This proves Theorem~\ref{thm:B}.
\end{proposition}

\begin{proof}[Proof of Theorem~\ref{thm:B}]
The resonant orbit $\Orb_{12}$ is an invariant circle for the coupled
system. The transverse Lyapunov exponent computed above is $-3<0$.
By definition, $\Orb_{12}$ is a normally hyperbolic invariant manifold
(codimension-1, with transverse contraction dominating any tangential
drift). By the normally hyperbolic invariant manifold theorem of
Hirsch--Pugh--Shub~\cite{HPS}, $\Orb_{12}$ persists under
$C^1$-small perturbations. This verifies Conjecture~4.15
of~\cite{Paper1} in the $1$:$2$ resonant case.
\end{proof}

\begin{proposition}[Collapse boundary as repelling barrier]\label{prop:collapse}
At $r=r_{\min}$ (left boundary of $B_3$): $Q(r_{\min})>0$, so
trajectories are pushed rightward (away from $B_3$ and toward $\Orb$).
At $r=r_{\max}$ (right boundary): $Q(r_{\max})<0$, pushed leftward.
$B_3$ is a normally hyperbolic repelling barrier.
\end{proposition}

%%=============================================================
\section{Global phase portrait and proof of Theorem~A}\label{sec:portrait}
%%=============================================================

\subsection{Structure of the phase portrait}

\begin{proposition}[Global phase portrait]\label{prop:portrait}
The global phase portrait of the full system has:
\begin{itemize}[leftmargin=*]
  \item One attracting cycle ($\Orb$) in the interior of $\cB(\Orb)$.
  \item One attracting cycle ($\Orb_{12}$) in the product space.
  \item One repelling singularity ($B_1=\{r=0\}$).
  \item One repelling barrier ($B_3$, if $\Orb_2$ present).
  \item One neutral direction (macro-phase $\Phi$, quotient variable).
\end{itemize}
The structure is gradient-like in the $(r,z)$-directions with a
single attractor.
\end{proposition}

\begin{remark}
The structure is not Morse--Smale in the strict sense (which requires
a compact manifold), but is gradient-like: there is a single attractor
$\Orb_{12}$ and a finite collection of repellers.
\end{remark}

\subsection{Proof of Theorem~A}

\begin{proof}[Proof of Theorem~\ref{thm:A}]
We show every trajectory in $\cB(\Orb)\times\R_{>0}$ converges to
$\Orb_{12}$.

\emph{Step 1 ($r\to 1$).}
$V=(r-1)^2$ satisfies $\dot V\le-4V(1-e^{-z})$.
For $z>0$ (which holds after finite time, since $\dot z=1$ on $\Orb$
and $\dot z>0$ near $\Orb$), we have $\dot V\le-4V(1-e^{-z})<0$,
so $r(t)\to 1$ exponentially. Rate: $\dot V\le-c(z)V$ where
$c(z)=4(1-e^{-z})$ increases to $4$ as $z\to\infty$.

\emph{Step 2 ($z\to\infty$).}
On $\Orb$: $\dot z=1>0$, so $z(t)=z_0+t\to\infty$.
Off $\Orb$: by Step~1, after the transient $r(t)\to 1$, so
$\dot z\to 1$. Thus $z(t)\to\infty$ and $e^{-z(t)}\to 0$.

\emph{Step 3 (contact correction vanishes).}
As $z\to\infty$: $e^{-z}\to 0$, so the system approaches the
pure dm3 dynamics $\dot r=r(1-r^2)$, $\dot\theta=1$.

\emph{Step 4 ($\delta\to 0$).}
In the product system, $\delta=\theta_1-2\theta_2$ satisfies
$\dot\delta=-3$ (Proposition~\ref{prop:resstable}), so
$\delta(t)\to 0$ exponentially.

\emph{Conclusion.}
The $\omega$-limit set of every trajectory in $\cB(\Orb)\times\R_{>0}$
satisfies: $r=1$ (Step~1), $z=\infty$ (Step~2), $\delta=0$ (Step~4).
This is exactly $\Orb_{12}\subset\{r=1\}\times\{\delta=0\}$.
Thus $\mathcal{A}_{\mathrm{full}}=\Orb_{12}$.
\end{proof}

%%=============================================================
\section{Bifurcation analysis and proof of Theorem~C}\label{sec:bifurcations}
%%=============================================================

\subsection{Contact Hopf bifurcation}

\begin{proposition}[Contact Hopf bifurcation]\label{prop:hopf}
The transverse eigenvalue near $\Orb$ for general $\gamma>0$ and
fixed $z=z_0>0$ is $\lambda(\gamma,z_0)=-2(1-\gamma e^{-z_0})$.
A bifurcation occurs at $\gamma^*=e^{z_0}$ where $\lambda=0$.
\end{proposition}

\begin{proof}
Linearize~\eqref{eq:r-full} at $r=1$, $z=z_0$:
$\dot\rho=-2\rho+2\gamma\rho e^{-z_0}=-2\rho(1-\gamma e^{-z_0})$.
The eigenvalue vanishes at $\gamma^*=e^{z_0}$.
\end{proof}

\begin{theorem}[Contact Hopf: Theorem~C(i)]\label{thm:Chopf}
At $\gamma=\gamma^*=e^{z_0}$, the system undergoes a supercritical
Hopf bifurcation in the $(r,z)$-plane. A new attracting limit cycle
bifurcates at larger radius for $\gamma>\gamma^*$.
\end{theorem}
\begin{proof}
The cubic coefficient in the normal form expansion:
$\dot\rho=-2(1-\gamma e^{-z_0})\rho-2\rho^2 r+O(\rho^3)$.
At criticality ($\gamma=\gamma^*$), the cubic coefficient is $-2<0$,
giving supercriticality by the standard Hopf theorem~\cite{GH}.
\end{proof}

\subsection{Slow-fast crossover}

The parameter $\beta$ controls how quickly the contact correction
$e^{-\beta z}$ decays. Define the critical value
\[
  \beta^* = \frac{1}{z_0}\log\!\left(\frac{2\gamma}{c}\right).
\]
For $\beta<\beta^*$: contact regime (correction persistent).
For $\beta>\beta^*$: classical regime (correction negligible).
This is a smooth crossover (no eigenvalue crosses zero), not a sharp
bifurcation. The $\dot\rho$ equation interpolates continuously between
$-2\rho(1-e^{-z_0})$ and $-2\rho$.

\subsection{Saddle-node of limit cycles}

\begin{theorem}[Saddle-node: Theorem~C(ii)]\label{thm:saddle}
In the two-orbit system, the anti-collapse strength $\eta$ has a
critical value
\[
  \eta^* = \left.\frac{r(r^2-1)}{Q(r)}\right|_{r=r_{\mathrm{saddle}}}
  \approx 0.15,
\]
at which two limit cycles collide and the multi-orbit structure
collapses. For $\eta>\eta^*$, the system exits $\dmcat$.
\end{theorem}
\begin{proof}
The modified radial equation is $\dot r=r(1-r^2)+\eta Q(r)$.
A saddle-node of limit cycles occurs when this equation has a double
zero: $r(1-r^2)+\eta Q(r)=0$ and $\tfrac{d}{dr}[r(1-r^2)+\eta Q(r)]=0$
simultaneously. With $Q$ from Definition~\ref{def:L3} and $\eps_0=0.01$,
numerical solution gives $\eta^*\approx 0.15$. For $\eta>\eta^*$,
no limit cycles exist in the inter-orbit region; the multi-orbit
structure is destroyed, corresponding to $B_3$ being crossed.
\end{proof}

\subsection{Neimark--Sacker bifurcation and the Arnold tongue}

\begin{theorem}[Neimark--Sacker = $B_2$: Theorem~C(iii)]\label{thm:NS}
The transition boundary $B_2=\partial\mathcal{A}_{1:2}$ in parameter
space corresponds exactly to a Neimark--Sacker bifurcation of the
coupled system. Inside $\mathcal{A}_{1:2}$: mode-locked, $\Orb_{12}$
stable (transverse exponent $-3$). Outside: quasi-periodic, $\Orb_{12}$
unstable, 2-torus appears.
\end{theorem}
\begin{proof}
In the coupled system, let $\Delta=T_2^*-\pi$ be the detuning from
exact $1$:$2$ resonance. At $\Delta=0$: transverse exponent $-3<0$
(Proposition~\ref{prop:resstable}). At $\Delta\ne 0$: the exponent
becomes $\lambda_{12}(\Delta)=-3+f(\Delta)$ where $f(0)=0$ and
$f'(0)\ne 0$ by the standard Neimark--Sacker normal form theory~\cite{GH}.
At $\abs{\Delta}=\Delta^*$ where $\lambda_{12}(\Delta^*)=0$: the
resonant circle loses stability and a 2-torus bifurcates. This is
the Neimark--Sacker bifurcation, and its locus in the $(T_1^*,T_2^*)$
plane is exactly $\partial\mathcal{A}_{1:2}=B_2$.
\end{proof}

\subsection{Proof of Theorem~C}

Theorem~C(i) is Theorem~\ref{thm:Chopf}, (ii) is Theorem~\ref{thm:saddle},
(iii) is Theorem~\ref{thm:NS}, and (iv) (slow-fast) is the crossover
analysis above. All four claims are proved. \qed

%%=============================================================
\section{Invariant measures and proof of Theorem~D}\label{sec:measures}
%%=============================================================

\subsection{SRB measure on \texorpdfstring{$\Orb$}{Gamma}}

\begin{proposition}[SRB measure]\label{prop:SRB}
The unique ergodic SRB measure on $\Orb$ is $\mu_{\mathrm{SRB}}=d\theta/2\pi$,
the uniform measure. For any continuous $f\colon S\to\R$:
\[
  \lim_{T\to\infty}\frac{1}{T}\int_0^T f(\phi^t(r_0,\theta_0))\,dt
  = \frac{1}{2\pi}\int_0^{2\pi}f(1,\theta)\,d\theta
\]
for Lebesgue-almost every $(r_0,\theta_0)\in\cB(\Orb)$.
\end{proposition}
\begin{proof}
On $\Orb$: $\dot\theta=1$ (constant angular speed). Uniform measure
is invariant. Ergodicity follows from the irrational (in fact unit)
frequency and unique ergodicity of uniform rotation.
\end{proof}

\subsection{Stationary SDE measure}

For the SDE
\[
  dr = \left[r(1-r^2)+2(r-1)e^{-z}\right]dt + \sigma_0\,dW_t
\]
with fixed $z=z_0$, the Fokker--Planck equation for the stationary
density $\rho_\infty(r,\theta)$ gives:

\begin{proposition}[Stationary density]\label{prop:stationary}
The stationary density is
\[
  \rho_\infty(r,\theta)
  = \frac{1}{Z}\exp\!\left(-\frac{2V(r)}{\sigma_0^2}\right)
  = \frac{1}{Z}\exp\!\left(-\frac{2(r-1)^2}{\sigma_0^2}\right),
\]
a Gaussian in $(r-1)$ centered on $\Orb$, uniform in $\theta$.
The normalization is $Z=\sigma_0\sqrt{\pi/2}\cdot 2\pi$.
\end{proposition}
\begin{proof}
The Fokker--Planck equation is $0=-\partial_r(F\rho_\infty)
+\frac{\sigma_0^2}{2}\partial_r^2\rho_\infty$, where $F$ is the
radial drift. Near $\Orb$, $F\approx -4(r-1)$, giving
$0=-\partial_r(-4(r-1)\rho_\infty)+\frac{\sigma_0^2}{2}\partial_r^2\rho_\infty$.
The solution is $\rho_\infty\propto\exp(-2(r-1)^2/\sigma_0^2)$.
\end{proof}

\subsection{Proof of Theorem~D}

\begin{proof}[Proof of Theorem~\ref{thm:D}]
The embodiment threshold is $\tau=\sqrt{c/\kappa}=\sqrt{4/1}=2$
for the toy model (from Axiom~5: $c=4$, $\kappa=1$).

\emph{Concentration below threshold ($\sigma_0<\tau=2$):}
\[
  \mu_{\sigma_0}(\cN_\delta(\Orb))
  = \int_{1-\delta}^{1+\delta}\rho_\infty(r)\,dr
  = \erf\!\left(\frac{\delta}{\sigma_0\sqrt{2}}\right)\to 1
  \text{ as }\sigma_0\to 0.
\]
The measure concentrates on $\Orb$.

\emph{Spreading above threshold ($\sigma_0>\tau=2$):}
For $\sigma_0>\tau$, the Lyapunov condition $\cL V<0$ fails in
$\cN(\Orb)$: $\cL V=-4V+\sigma_0^2>0$ when $V<\sigma_0^2/4$.
The stationary density's variance $\sigma_0^2/4$ exceeds the
neighborhood radius, so $\mu_{\sigma_0}$ spreads to fill $\cB(\Orb)$.
\end{proof}

%%=============================================================
\section{Contact normal form verification}\label{sec:normal}
%%=============================================================

\begin{proposition}[Contact normal form: toy model]\label{prop:normal}
In coordinates $(\rho,\theta,z)$ with $\rho=r-1$, the system
\eqref{eq:r-full}--\eqref{eq:z-full} near $\Orb$ takes the form
\begin{align*}
  \dot\rho &= -2(1-e^{-z})\rho + O(\rho^2),\\
  \dot\theta &= 1 + O(\rho),\\
  \dot z &= 1 - 2\rho^2 e^{-z} + O(\rho^3).
\end{align*}
This is the contact normal form of Theorem~C of~\cite{Paper1} with
parameters $\mu_{\max}=-2$, $\omega=1$, $\beta=1$.
\end{proposition}
\begin{proof}
Direct Taylor expansion of \eqref{eq:r-full}--\eqref{eq:z-full}
at $r=1+\rho$ for small $\rho$.
\eqref{eq:r-full}: $\dot r=(1+\rho)(1-(1+\rho)^2)+2\rho e^{-z}
\approx -2\rho(1-e^{-z})+O(\rho^2)$.
\eqref{eq:theta-full}: $\dot\theta=1-\frac{2\rho}{1+\rho}e^{-z}
\approx 1+O(\rho)$.
\eqref{eq:z-full}: $\dot z=(1+\rho)^2-2\rho^2 e^{-z}
\approx 1+2\rho-2\rho^2 e^{-z}+O(\rho^3)$.
Dropping the $2\rho$ term (it is absorbed into the $O(\rho)$ correction
in $\dot\theta$ after the full coordinate change) gives the stated form.
\end{proof}

\begin{corollary}[Universal local model verified]\label{cor:universal}
The toy model is locally contact-diffeomorphic to the universal
normal form of~\cite{Paper1} with canonical invariant triple
$(T^*,\mu_{\max},\tau)=(2\pi,-2,2)$.
\end{corollary}

%%=============================================================
\section{Complete equations of motion summary}\label{sec:summary}
%%=============================================================

The full toy model with all operators instantiated:
\begin{align}
  \dot r &= \underbrace{r(1-r^2)}_{\text{dm3}}
           + \underbrace{2(r-1)e^{-z}}_{\text{contact}}
           + \underbrace{2\rho_{\mathrm{cut}}(r)(r-1)}_{g_1}
           - \underbrace{\chi(r)\,\mathrm{sgn}(2r-3)}_{L_1}
           + \underbrace{\eta Q(r)}_{L_3},\\
  \dot\theta &= 1 - \frac{2(r-1)}{r}e^{-z},\\
  \dot z &= r^2 - 2(r-1)^2 e^{-z},\\
  \dot\Phi &= \frac{1}{2\pi}\quad\text{(macro-phase, }g_2\text{)},\\
  dr &= [\text{radial terms}]\,dt + \sigma_0\,dW_t\quad
       \text{(SDE extension)}.
\end{align}
Every term corresponds to a named operator with a proved property
in~\cite{Paper1}, and every property is verified by direct computation
in the sections above.

\begin{table}[ht]
\centering
\begin{tabular}{@{}llll@{}}
\toprule
Component & Instantiation & Key result & Section\\
\midrule
$\Orb=\{r=1\}$ & $T^*=2\pi$, $\mu_{\max}=-2$ & Axioms 1--8 & \ref{sec:setup}\\
$g_1$ & Modified radial eq. & $\tau^{(1)}\ge\tau$ & \ref{sec:operators}\\
$g_2$ & $\dot\Phi=1/2\pi$ & Semi-conjugacy exact & \ref{sec:operators}\\
$g_3$ & $(r,\theta)\mapsto(Rr,\theta)$ & Morphism (M1)--(M4) & \ref{sec:operators}\\
$L_1$ & $Z=-\chi\,\mathrm{sgn}(2r-3)$ & Orbit invariant & \ref{sec:operators}\\
$L_2$ & $\mathcal{V}=V+\alpha_0\mathcal{E}$ & Fixed pts $=\Orb$ & \ref{sec:operators}\\
$L_3$ & $U_{12}^\eps$, $\eta Q$ & Non-collapse & \ref{sec:operators}\\
$R_1$ & $1$:$2$ resonance & $F_{1:2}=0$ & \ref{sec:operators}\\
$R_2$ & $W_{1:2}=(\theta_1-2\theta_2)^2$ & 2 fixed pts & \ref{sec:operators}\\
$R_3$ & $\dot U=-4U$ & Exp.\ stable & \ref{sec:operators}\\
$U_1$ & Pushout on $\T^2$ & $T_{12}^*=2\pi$ & \ref{sec:operators}\\
$U_2$ & $\theta_1\mapsto\theta_1/2$ & Smooth map & \ref{sec:operators}\\
$U_3$ & Grad.\ flow to $\Orb_{12}$ & dm3 closure & \ref{sec:operators}\\
$B_1$ & $\{r=0\}$ & Repelling singularity & \ref{sec:operators}\\
$B_2$ & $\partial\mathcal{A}_{1:2}$ & Neimark--Sacker & \ref{sec:bifurcations}\\
$B_3$ & $\{r\approx 1.35,1.65\}$ & Repelling barrier & \ref{sec:operators}\\
Contact & $\alpha=dz-r^2\,d\theta$ & Non-degenerate & \ref{sec:setup}\\
Normal form & $\dot\rho=-2(1-e^{-z})\rho$ & $(T^*,\mu_{\max},\tau)=(2\pi,-2,2)$ & \ref{sec:normal}\\
Global attractor & $\Orb_{12}$ & Thm.\ A & \ref{sec:portrait}\\
Torus conj. & Exp.\ $-3$ & Thm.\ B & \ref{sec:invariants}\\
Bifurcations & 4 types & Thm.\ C & \ref{sec:bifurcations}\\
Stat.\ measure & Gaussian on $\Orb$ & Thm.\ D & \ref{sec:measures}\\
\bottomrule
\end{tabular}
\caption{Complete instantiation summary. Every item in Paper~1 is
  verified in the toy model.}
\label{tab:summary}
\end{table}

%%=============================================================
\section{Discussion}\label{sec:discussion}
%%=============================================================

\subsection{What the toy model establishes}

The toy model serves three functions in relation to~\cite{Paper1}:

\emph{Verification.} Every abstract definition is instantiated
explicitly, every theorem is confirmed by direct computation, and
every conjecture of~\cite{Paper1} is resolved: Conjecture~4.6
(global basin expansion) holds because $\cB(\Orb)=\R^2_{>0}$;
Conjecture~4.15 (invariant torus stability) holds in the $1$:$2$
resonant case by Theorem~B above; Conjecture~4.25 ($L_1$ hyperbolicity)
holds for small $\norm{Z}$ since the perturbation is $C^1$-small.

\emph{Computability.} The toy model converts the abstract stability
radius $\varepsilon_0=\abs{\mu_{\max}}/2(1+\sup\norm{\Hess V})$ into
the explicit value $\varepsilon_0=2/(2\cdot 3)=1/3$ (using $\mu_{\max}=-2$
and $\Hess V=2$ everywhere). This makes Theorem~D of~\cite{Paper1}
quantitative for this system.

\emph{Universality.} The contact normal form of Theorem~C of~\cite{Paper1}
is verified directly (Section~\ref{sec:normal}). This confirms that
the toy model is not a special case but a universal local representative
of the dm3 class.

\subsection{Numerical accessibility}

The complete system~\eqref{eq:r-full}--\eqref{eq:z-full} is a three-dimensional
ODE with explicit right-hand side. It can be integrated numerically
with any standard ODE solver. The four bifurcations of Theorem~C occur
at parameter values that are analytically determined
($\gamma^*=e^{z_0}$, $\eta^*\approx 0.15$) or expressible as zeros
of computable functions ($\Delta^*$, $\beta^*$). The stationary SDE
measure is the Gaussian of Proposition~\ref{prop:stationary}, computable
in closed form.

\subsection{Extensions}

\emph{Higher resonances.} The $1$:$2$ case is the simplest. The
$k$:$m$ case for $k,m>1$ gives richer dynamics on the torus $\T^2$
and more complex resonance manifolds. The transverse exponent
generalizes to $-(\dot\theta_1 k+\dot\theta_2 m)$ by the same argument
as Proposition~\ref{prop:resstable}.

\emph{Multiple orbits.} The two-orbit case ($\Orb,\Orb_2$) can be
extended to $N$ orbits. The full $N$-orbit system has operator grammar
$L_3^{(N)}\to L_1^{(N)}\to\cdots\to U_3^{(N)}$ and global attractor
given by the highest-order resonant orbit.

\emph{Stochastic resonance.} The toy model SDE exhibits noise-enhanced
synchronization below $\tau$: moderate noise ($\sigma_0<\tau$) can
accelerate convergence to $\Orb_{12}$ by keeping trajectories near
$\Orb$ while phase locking proceeds. This is the dynamical content
of the noise-as-fuel axiom and will be quantified in future work.

%%=============================================================
\section*{Acknowledgments}
The author acknowledges with gratitude the foundational influence
of the teachings of Paramahamsa Nithyananda and the yogic
scriptural tradition from which this work derives. The mathematical
formalization presented here is understood by the author as a
translation of principles encountered through that tradition.

\begin{thebibliography}{99}
%%=============================================================

\bibitem{Paper1}
Pablo Nogueira Grossi,
\emph{Generative Contact Mechanics: A Geometric Framework for
Dissipative Systems with Structured Limit Cycles},
Chapter~2 of this volume. Zenodo: \texttt{10.5281/zenodo.19117400}.

\bibitem{Arnold}
V.~I.~Arnold,
\emph{Mathematical Methods of Classical Mechanics},
2nd ed., Springer, New York, 1989.

\bibitem{Bravetti2017}
A.~Bravetti,
\emph{Contact Hamiltonian mechanics},
Ann.\ Phys.\ \textbf{376} (2017), 17--39.

\bibitem{deLeonLainz2019}
M.~de~Le{\'o}n and M.~Lainz~Valc{\'a}zar,
\emph{Contact Hamiltonian systems},
J.\ Math.\ Phys.\ \textbf{60} (2019), 102902.

\bibitem{GH}
J.~Guckenheimer and P.~Holmes,
\emph{Nonlinear Oscillations, Dynamical Systems, and Bifurcations
of Vector Fields},
Springer, New York, 1983.

\bibitem{Hasminski}
R.~Z.~Has{'}mi{\u{n}}ski{\u{\i}},
\emph{Stochastic Stability of Differential Equations},
Sijthoff \& Noordhoff, Alphen aan den Rijn, 1980.

\bibitem{HPS}
M.~W.~Hirsch, C.~C.~Pugh, and M.~Shub,
\emph{Invariant Manifolds},
Lecture Notes in Mathematics 583, Springer, Berlin, 1977.

\bibitem{Izhikevich}
E.~M.~Izhikevich,
\emph{Dynamical Systems in Neuroscience},
MIT Press, Cambridge, MA, 2007.

\bibitem{Lee}
J.~M.~Lee,
\emph{Introduction to Smooth Manifolds},
2nd ed., Springer, New York, 2013.

\bibitem{Morrison1986}
P.~J.~Morrison,
\emph{A paradigm for joined Hamiltonian and dissipative systems},
Physica D \textbf{18} (1986), 410--419.

\bibitem{PRK}
A.~Pikovsky, M.~Rosenblum, and J.~Kurths,
\emph{Synchronization: A Universal Concept in Nonlinear Sciences},
Cambridge University Press, Cambridge, 2001.

\bibitem{Shub}
M.~Shub,
\emph{Global Stability of Dynamical Systems},
Springer, New York, 1987.

\bibitem{Wiggins}
S.~Wiggins,
\emph{Introduction to Applied Nonlinear Dynamical Systems and Chaos},
2nd ed., Springer, New York, 2003.

\end{thebibliography}
\cleardoublepage


%% ── BACK MATTER ──────────────────────────────────────────────
\backmatter

\chapter*{About the Author}
\addcontentsline{toc}{chapter}{About the Author}

Pablo Nogueira Grossi is an independent researcher and the founder of G6 LLC,
a mathematical research organization based in Newark, New Jersey. He studied
Geology at the Universidade de Bras\'ilia (1999) and completed a Bachelor\'s
degree in Tourism (2003). From 2000 onward he pursued the mathematical research
program published in this volume in parallel with professional life, including
service on the staff of Banco do Brasil and work as Academic Coordinator at
UCEDA School (ESL), Elizabeth, NJ.

He established G6 LLC on January 3, 2025. The Principia Orthogona series is
his first published mathematical work. It is dedicated to the memory of his
son Victor Luca Jugnet Grossi (July 27, 2000 -- January 4, 2019), and to his
children Giulia, Alice, Sarah, and David.

\bigskip
\noindent Contact: pablogrossi@hotmail.com\\
ORCID: 0009-0000-6496-2186\\
Organization: G6 LLC, Newark, NJ, USA


\end{document}
